% Generated from the modular source by build.py; edit the modular files.
\documentclass[11pt,a4paper]{article}
\usepackage[T1]{fontenc}
\usepackage[utf8]{inputenc}
\usepackage{lmodern}
\usepackage[a4paper,margin=22mm,headheight=14pt,headsep=8mm]{geometry}
\usepackage{amsmath,amssymb,amsthm,mathtools,mathrsfs}
\usepackage{booktabs,array,longtable,enumitem}
\usepackage{microtype}
\usepackage[dvipsnames]{xcolor}
\usepackage{fancyhdr}
\usepackage{hyperref}
\hypersetup{colorlinks=true,linkcolor=MidnightBlue,citecolor=MidnightBlue,
 urlcolor=MidnightBlue,pdftitle={Exact Identities for Polylogarithms, Harmonic Zeta Functions, and Stieltjes Finite Parts},
 pdfauthor={Research continuation for the ProveIt project},pdfsubject={Exact identities and rigorous finite-part calculus}}
\usepackage{bookmark}
\numberwithin{equation}{section}
\newtheorem{theorem}{Theorem}[section]
\newtheorem{proposition}[theorem]{Proposition}
\newtheorem{lemma}[theorem]{Lemma}
\newtheorem{corollary}[theorem]{Corollary}
\theoremstyle{definition}
\newtheorem{definition}[theorem]{Definition}
\theoremstyle{remark}
\newtheorem{remark}[theorem]{Remark}
\DeclareMathOperator{\Li}{Li}
\DeclareMathOperator{\FP}{FP}
\DeclareMathOperator{\CT}{CT}
\setlist{topsep=4pt,itemsep=3pt,parsep=0pt}
\setlength{\parskip}{2pt}
\setlength{\emergencystretch}{2em}
\allowdisplaybreaks[2]
\pagestyle{fancy}
\fancyhf{}
\fancyhead[L]{\small Exact special-function identities}
\fancyhead[R]{\small ProveIt research continuation}
\fancyfoot[C]{\thepage}
\renewcommand{\headrulewidth}{0.3pt}
\title{\vspace{-7mm}\textbf{Exact Identities for Polylogarithms,\\
Harmonic Zeta Functions, and\\Stieltjes Finite Parts}\\[4mm]
\large Colliding grids, nonlinear coordinates, fully shifted sums,\\
and trilogarithmic antiderivatives}
\author{Research continuation prepared for Vladimir Reshetnikov\\
\small The ProveIt project}
\date{10 October 2026}
\begin{document}
\maketitle
\thispagestyle{empty}
\begin{abstract}
We continue the exact-identity program in the supplied ProveIt manuscript
and its five current incoming polylogarithm reports. We obtain a
holomorphic coefficient formula for every colliding unequal-dilation
product of generalized Stieltjes functions and their argument derivatives,
including a convergent collision expansion and an odd-order vanishing
family. A finite Taylor-coefficient law determines every nonlinear
coordinate correction, its composition rule, and its antiderivative
endpoint constant. Three separated unequal frequencies reduce to a finite
family of colored Tornheim kernels; complete local subtraction then gives
all Stieltjes and argument derivatives. For fully shifted height-one
harmonic sums, a Gamma-normalized difference is entire in the order, and
all its nonpositive-integer values are finite Stirling transforms. This
yields every Laurent principal part and finite part, with exact pole
orders and Gamma antiderivatives. Finally, the weighted arctangent family
left unevaluated in the canonical chapter is reduced to two values of one
real trilogarithmic function, together with ordinary primitives and a
double-polylogarithm integration ladder. Proofs distinguish coordinate
finite parts from meromorphic coefficient extraction. The report records
two pending source corrections, preserves the unresolved status of the
short $S_6$ and current $S_8$ conjectures, and proposes precise next
identity problems. Numerical and symbolic audits accompany, but do not
replace, the proofs.
\end{abstract}
\vfill
\noindent\textbf{Source baseline.} Repository commit
\texttt{445f754610e1377939235794de84ed9075fc09f5}.\\
\textbf{Status.} Analytic research manuscript with executable audits;
no Lean formalization or exhaustive historical-priority claim.
\clearpage
\begingroup
\small
\setlength{\parskip}{0pt}
\tableofcontents
\endgroup
\clearpage

% BEGIN sections/introduction.tex
\section{Scope, results, and proof status}
\label{xid:sec:introduction}

This report addresses exact identities and the analytic mechanisms that
make them valid. The source is the unified manuscript at
\texttt{Analysis/Polylogarithms/docs/manuscript}, together with the five
polylogarithm archives present in \texttt{docs/incoming} at the commit
printed on the title page \cite{ProveIt,IncomingGauss,IncomingMellin,
IncomingNested,IncomingContinuation,IncomingIdentities}. All 85 retrieved
file payloads were matched to their Git blob identifiers before the
research comparison. The accompanying provenance manifest fixes precisely
what was inspected; later changes to the repository are outside this
comparison.

Several questions in those reports concern the point where an exact
formula is known in a regular domain but its endpoint extension has not
yet been justified. Typical examples are a collision of two rational
singularity grids, replacing an affine endpoint coordinate by a nonlinear
one, and taking all argument derivatives of a three-factor integral.
We resolve these questions under explicit hypotheses. Another direction
is a fully shifted nested harmonic family, where changing every
denominator produces an entire correction that the outer-shift theory
does not contain. A final direction returns to the canonical weighted
arctangent integral and gives its complete real evaluation.

\subsection{The results at a glance}

\begin{enumerate}[label=\textbf{\arabic*.},leftmargin=*]
\item \textbf{Every colliding two-grid Stieltjes product.}
Theorem~\ref{xid:cld:master} gives a finite holomorphic generator for all
Stieltjes indices, all argument derivative orders of positive total order,
and all positive integer dilations. Section~\ref{xid:cld:section} also treats
total order zero. The coefficients close in ordinary zeta derivatives,
harmonic numbers, and logarithms. The regular constant in the convergent
separated-shift collision expansion is exactly the independently defined
colliding finite part. Corollary~\ref{xid:cld:diagonal-collapse} gives an
additional all-index odd-order cancellation.

\item \textbf{Every nonlinear endpoint correction.}
Theorem~\ref{xid:thm:coordinate-monomial} gives the complete change of a
logarithmic-pole finite part from a finite jet of the coordinate.
Theorems~\ref{xid:thm:cocycle} and \ref{xid:thm:stieltjes-nonlinear} give its
composition and Stieltjes forms. A sharp disappearance threshold holds
for coordinates tangent to the identity, and
Theorem~\ref{xid:thm:nonlinear-primitive} gives exact endpoint constants for
ordinary primitives evaluated by a coordinate finite part.

\item \textbf{Three unequal separated frequencies.}
Theorem~\ref{xid:thm:triple-spectral} gives a finite sign-cone and character
formula involving $2(p_1+p_2+p_3)$ weighted colored Tornheim kernels.
Each weighted kernel is itself a finite sum of ordinary colored kernels.
Theorem~\ref{xid:thm:triple-FP} performs complete ambient-coordinate
subtraction for every Stieltjes index and argument derivative.

\item \textbf{An entire difference for fully shifted harmonic sums.}
Theorem~\ref{xid:hsh:thm:entire} isolates a universal Gamma-normalized
singular part and evaluates the entire remainder at every nonpositive
integer by a finite Stirling sum. Theorem~\ref{xid:hsh:thm:Laurent} gives
the complete Laurent principal part and finite part at all depths.
The pole order drops at negative even integers for a proved Bernoulli
reason. A two-shift bridge and explicit Gamma antiderivatives follow.

\item \textbf{Weighted arctangent and inverse-hyperbolic identities.}
Theorem~\ref{xid:wat:thm:main} evaluates the canonical weighted family by a
hyperbolic addition formula for one even trilogarithmic function.
Theorem~\ref{xid:wat:thm:primitive} supplies a normalized ordinary primitive
at every endpoint in $(0,\pi)$, with principal branches justified.
Theorem~\ref{xid:wat:thm:ladder} gives every Euler integration order in
depth-two polylogarithms. A trigonometric companion evaluates a weighted
inverse-hyperbolic family, including its endpoint value.
\end{enumerate}

These are proved extensions relative to the inspected source material.
They are not statements that every displayed identity is historically
new. In particular, the Hurwitz Fourier product, Gamma multiplication,
Gauss connection formula, and elementary logarithmic primitives are
classical inputs \cite{EspinosaMoll1,EspinosaMoll2,DLMFhurwitz,
hsh:DLMFGauss,wat:DavidH}. The depth-one harmonic Hurwitz-zeta literature
already treats continuation and Laurent coefficients
\cite{hsh:Kargin2023,hsh:KaraCan2026}; the limited access to the newest
preprint is disclosed in Section~\ref{xid:hsh:sec:main}.

\subsection{What the proofs have in common}

A local term of the form $c(z)t^{z-1}$ explains much of the endpoint
calculus. Its meromorphic integral over $(0,h)$ is $c(z)h^z/z$.
The generating function of its coefficientwise, fixed-coordinate finite
parts is instead
\[
 \frac{c(z)(h^z-1)}z.
\]
The difference is the complete expression $c(z)/z$. Retaining only its
residue leaves incorrect regular coefficients whenever $c$ depends on
$z$. Nonlinear coordinates, dilations, and Taylor coefficients of other
factors all produce precisely this dependence. The local subtraction
proofs below retain it before taking any coefficient.

The fully shifted harmonic theorem uses the complementary mechanism.
Its Gauss connection formula splits the Mellin kernel into a universal
singular block and an analytic block. Division by $\Gamma(s)$ cancels
the analytic block's simple Mellin poles. At $s=-m$ the resulting entire
function is determined by one Taylor coefficient, which turns out to be
a finite Stirling polynomial. Thus exact local data control both the
distributional correction and the entire harmonic correction, although
the two operations are analytically different.

\subsection{Conjectures not settled by this report}

The source's $S_4$ identity is already proved by its octahedral rational
certificate. The short $S_6$ candidate and the current, replacement $S_8$
candidate remain conjectures. The rejected earlier $S_8$ vector remains
rejected. The numerically supported higher golden ladder vectors are
also not proved here. The present exact results do not imply those
reductions, and their numerical evidence has not been used as a proof.
Section~\ref{xid:sec:audit} records exact source labels and the integration
consequences.

% END sections/introduction.tex


% BEGIN sections/preliminaries.tex
\section{Conventions and analytic preliminaries}
\label{xid:sec:preliminaries}

\subsection{Order derivatives, argument derivatives, and harmonic numbers}

For $\Re a>0$, the Hurwitz zeta function is initially
\[
 \zeta(s,a)=\sum_{n=0}^{\infty}(n+a)^{-s},\qquad\Re s>1,
\]
with the principal logarithm, and then has its usual meromorphic
continuation. We use
\begin{equation}\label{xid:eq:Stieltjes-convention}
 \zeta(s,a)=\frac1{s-1}
  +\sum_{n\ge0}\frac{(-1)^n}{n!}\gamma_n(a)(s-1)^n,
 \qquad \gamma_n=\gamma_n(1),\quad\gamma_0=\gamma.
\end{equation}
Thus the generator near $z=0$ is
\begin{equation}\label{xid:eq:G-convention}
 \mathscr G(z,a)=\zeta(1-z,a)+\frac1z
       =\sum_{n\ge0}\gamma_n(a)\frac{z^n}{n!},
 \qquad\gamma_0(a)=-\psi(a).
\end{equation}
A superscript on $\zeta$, as in $\zeta^{(j)}(s,a)$, means
$\partial_s^j$; a superscript on $\gamma_n$ or $\psi$ means an
argument derivative. The notation $D_y$ is used for a distributional
derivative. These three operations must be distinguished.
The standard identities used below are
\begin{align}
 \zeta(s,a)&=a^{-s}+\zeta(s,a+1),\label{xid:eq:Hurwitz-shift}\\
 \partial_a^r\zeta(s,a)&=(-1)^r(s)_r\zeta(s+r,a),\label{xid:eq:Hurwitz-arg}\\
 \psi^{(r)}(a)&=(-1)^{r+1}r!\zeta(r+1,a)\quad(r\ge1).
 \label{xid:eq:psi-Hurwitz}
\end{align}
They fix the normalizations and agree with DLMF
\cite{DLMFhurwitz,hsh:DLMFPolygamma}.

For later coefficient algebra put
\[
 H_N^{(s)}(a)=\sum_{j=0}^{N-1}(a+j)^{-s}
             =\zeta(s,a)-\zeta(s,N+a),
 \qquad H_N^{(s)}=H_N^{(s)}(1).
\]
This finite sum is entire in $s$; every order derivative is obtained by
inserting $(-\log(a+j))^k$. For integer $s\ge2$, it is also a difference
of polygamma values by \eqref{xid:eq:psi-Hurwitz}. We set $H_0^{(s)}=0$.
The elementary symmetric functions of $1,1/2,\ldots,1/r$ are finite
polynomials in $H_r,H_r^{(2)},\ldots$ by Newton's identities.
Coefficient extraction $[z^nt^j]$ always denotes the ordinary Taylor
coefficient; factorials are printed explicitly.

For example, \eqref{xid:eq:Hurwitz-arg} gives the useful classical bridge
\begin{equation}\label{xid:eq:Stieltjes-arg-zeta}
 \gamma_n^{(r)}(a)=(-1)^{n+r}r!
 \sum_{h=0}^{\min(n,r)}\frac{n!}{(n-h)!}
 e_h(1,1/2,\ldots,1/r)\,
       \zeta^{(n-h)}(r+1,a),\qquad r\ge1.
\end{equation}
Indeed, take the coefficient of $z^n$ in
$(-1)^r(1-z)_r\zeta(r+1-z,a)$. An ordinary primitive is
\begin{equation}\label{xid:eq:ordinary-Stieltjes-primitive}
 \int\gamma_n(a)\,da
       =\frac{(-1)^{n+1}}{n+1}\zeta^{(n+1)}(0,a)+C.
\end{equation}
One proof takes coefficients in
$\partial_a\zeta(-z,a)=z\zeta(1-z,a)$ after removing the constant
spectral pole. At $n=0$ this agrees with $-\log\Gamma(a)+C$.
These identities are baseline transformations, not claims of new
results; the endpoint constants after nonlinear substitution require
the additional analysis of Section~\ref{xid:sec:nonlinear}.

\subsection{The coordinate finite part}

Suppose a function has finitely many singular points and, to the right
of each point, a finite sum of nonintegrable terms
$t^{-q}(\log t)^n$ plus a locally integrable remainder. Integrate
outside right intervals of lengths $\varepsilon_j$ and expand in these
cutoffs. The symbol $\FP_x$ means the coefficient of
$\prod_j\varepsilon_j^0(\log\varepsilon_j)^0$ in this expansion,
with $t=x-x_j$ the original coordinate. Independent cutoffs are
permitted. A common cutoff gives the same constant because the finitely
many endpoint contributions add. No left cutoff is needed for a
fractional-part Hurwitz singularity: its left limit is the regular
branch at argument one.

For distributions the same prescription is applied after pairing with
a smooth test function and subtracting its Taylor polynomial. It is
local and linear. Products in this report are finite parts of
\emph{pointwise products} with explicitly subtracted local expansions.
No undefined multiplication of colliding singular distributions is
invoked. In contrast, a coefficient of a meromorphic family is a
spectral operation; the two operations coincide only after the required
complete local pole terms have been accounted for.

The periodic Hurwitz distribution $Z(1-z)$ is defined by the ordinary
locally integrable function for $\Re z>0$ and meromorphically continued
by local Taylor subtraction. At an integer, in a signed local
coordinate $y$, its decomposition is
\[
 Z(1-z)=y_+^{z-1}+\zeta(1-z,1+y).
\]
Consequently
\begin{equation}\label{xid:eq:canonical-periodic}
 Z(1-z)=\frac{\delta_0-1}{z}
       +\sum_{n\ge0}G_n\frac{z^n}{n!},
 \qquad G_n=\FP_y\gamma_n(\{y\}).
\end{equation}
Here $1$ means the constant periodic distribution. This construction
also fixes the normalization of every $D_y^rG_n$ used below.

\subsection{Fourier normalization and the two-frequency kernel}

Write $e(t)=e^{2\pi it}$. For $\Re z>1$, the absolutely convergent
Hurwitz Fourier series in our convention is
\begin{equation}\label{xid:eq:Hurwitz-Fourier}
 \zeta(1-z,x)=\frac{\Gamma(z)}{(2\pi)^z}
 \sum_{n\ne0}\frac{e^{-i\pi z\operatorname{sgn}(n)/2}}{|n|^z}e(nx).
\end{equation}
For positive integers $p,q$ set $d=\gcd(p,q)$,
$P=p/d$, $Q=q/d$, and $\tau=Qa-Pb$. Orthogonality leaves
frequencies $Q\ell,-P\ell$ and therefore gives
\begin{align}
 C(z,w;p,q;a,b)
 &: =\int_0^1\zeta(1-z,\{px+a\})
                  \zeta(1-w,\{qx+b\})\,dx\notag\\
 &=\frac{\Gamma(z)\Gamma(w)Q^{-z}P^{-w}}{(2\pi)^{z+w}}
 \left\{e^{-i\pi(z-w)/2}\Li_{z+w}(e(\tau))
       +e^{i\pi(z-w)/2}\Li_{z+w}(e(-\tau))\right\}.
 \label{xid:eq:pair-spectral}
\end{align}
Initially $\Re z,\Re w>1$; subsequently the formula means
meromorphic continuation. For $\tau\in\mathbb Z$ the two
polylogarithms are $\zeta(z+w)$. The unit-frequency coincident case
is the classical Espinosa--Moll product \cite[Theorem 3.1]{EspinosaMoll1}.
The shifted and dilated versions in the incoming reports have this
normalization. Formula~\eqref{xid:eq:pair-spectral} makes the pair kernels
needed in Section~\ref{xid:sec:triple-dilation} explicit within this report.

The Riemann and Hurwitz functional equations, Gamma reflection and
multiplication, and polylogarithm order continuation are used in their
standard forms \cite{DLMFhurwitz,wat:DLMF}. Each use of a singular
parameter value below is licensed by a local continuation proof before
its Taylor or Laurent coefficient is taken.

% END sections/preliminaries.tex


% BEGIN sections/collision.tex
% Research contribution prepared independently for the root article.
% Needs amsmath, amssymb, amsthm.  All labels have prefix cld:.
% Suggested references: EspinosaMoll2002, DLMFhurwitz, IncomingMellinDilation,
% IncomingStieltjesClosure.  Classical spectral kernels are credited inputs.

\section{Colliding unequal dilations: a complete finite-part calculus}
\label{xid:cld:section}

The separated-grid theorem in the incoming Mellin--dilation report
\cite{IncomingMellin}
explicitly excludes $Pa\in\mathbb Z$.  At these shifts its partition of
unity no longer defines a product of the two singular distributions.
This section treats precisely that omitted case.  The construction is
the constant term of the pointwise product in a fixed ambient coordinate;
it neither assumes that a product of colliding distributions exists nor
identifies a bare spectral Laurent coefficient with that constant term.

Let $p,q$ be positive integers and put
\[
 d=\gcd(p,q),\qquad P=p/d,\qquad Q=q/d,\qquad
 A=\log P,\quad B=\log Q,\quad
 K=\log\operatorname{lcm}(p,q).
\]
For $m,n,r,k\ge0$, write $N=r+k$ and define
\begin{equation}
 J_{mn;r,k}^{p,q}
 =\operatorname{FP}_x\int_0^1
       \gamma_m^{(r)}(\{px\})\gamma_n^{(k)}(\{qx\})\,dx.
 \label{xid:cld:def:J}
\end{equation}
The derivatives are in the special-function argument.  At every member
of the union of the two singular grids, delete the right interval of
ambient length $\epsilon$ and retain the constant term as
$\epsilon\downarrow0$.  Equivalently, take separate constant terms in
independent ambient cutoff variables at the finitely many points.

Every colliding shift $a$ gives the same value as \eqref{xid:cld:def:J}.
Indeed, $Pa=h\in\mathbb Z$, and a solution of
$Qj\equiv-h\pmod P$ makes $x_0=j/p$ satisfy
$px_0,qx_0+a\in\mathbb Z$.  Translation by $x_0$ carries the two
fractional-part arguments to $px$ and $qx$ and preserves ambient lengths.
Thus no collision cases are lost by writing the common shift as zero.

\subsection{The complete local subtraction polynomial}

Use the holomorphic Stieltjes generating function
\[
 \mathscr G(z,x)=\zeta(1-z,x)+\frac1z
       =\sum_{m\ge0}\gamma_m(x)\frac{z^m}{m!},
 \qquad
 p_r(z)=\frac{(1-z)_r}{r!}
       =\prod_{j=1}^r\left(1-\frac zj\right).
\]
For $r=0$ the empty product is one.  Set
\begin{equation}
 S_{m,r}(L)=m![z^m]\,p_r(z)e^{zL}
 =\sum_{j=0}^{\min(m,r)}(-1)^j\frac{m!}{(m-j)!}
     e_j(1,1/2,\ldots,1/r)L^{m-j}.
 \label{xid:cld:S-polynomial}
\end{equation}
The Hurwitz shift identity gives the exact local decomposition
\begin{equation}
 \gamma_m^{(r)}(y)
 =(-1)^r r!y^{-r-1}S_{m,r}(\log y)
      +\gamma_m^{(r)}(1+y).
 \label{xid:cld:local-decomposition}
\end{equation}

Here is an explicit subtraction rule that covers every grid configuration.
At a point $b$ of the first grid put $t=x-b>0$ and
\[
 U_{p,m,r}(t)=(-1)^r r!p^{-r-1}t^{-r-1}
                     S_{m,r}(\log p+\log t).
\]
If $b$ is not on the second grid, subtract
\[
 U_{p,m,r}(t)\sum_{j=0}^r
       \frac{q^jt^j}{j!}\gamma_n^{(k+j)}(\{qb\}).
\]
At a point of the second grid only, use the analogous expression with
the two factors exchanged.  At a common point subtract the sum of
\begin{align}
 &U_{p,m,r}(t)U_{q,n,k}(t),\\
 &U_{p,m,r}(t)\sum_{j=0}^r
       \frac{q^jt^j}{j!}\gamma_n^{(k+j)}(1),\\
 &U_{q,n,k}(t)\sum_{j=0}^k
       \frac{p^jt^j}{j!}\gamma_m^{(r+j)}(1).
 \label{xid:cld:common-subtraction}
\end{align}
The first line has power $t^{-N-2}$ and a polynomial in $\log t$.
The other lines contain the powers $t^{-r-1},\ldots,t^{-1}$ and
$t^{-k-1},\ldots,t^{-1}$.  These are all nonintegrable terms.

Taylor's theorem and \eqref{xid:cld:local-decomposition} show that the
remaining integrand is $O(1+|\log t|^{m+n})$.  It is therefore integrable.
Adding the elementary finite-part primitives of the subtracted terms
proves existence of \eqref{xid:cld:def:J}.  For clarity, the primitive
normalization is
\[
 \operatorname{FP}\int_0^\ell t^{-1}\log^j t\,dt
      =\frac{\log^{j+1}\ell}{j+1},
\]
and, for $h\ge2$,
\[
 \operatorname{FP}\int_0^\ell t^{-h}\log^j t\,dt
 =\left.\partial_\alpha^j
       \frac{\ell^{\alpha-h+1}}{\alpha-h+1}\right|_{\alpha=0}.
\]
No finite local constant is omitted or added.

\subsection{A holomorphic master kernel for every positive derivative order}

For $M\ge1$ define
\begin{align}
 W_M(z)&=(1-z)_M\zeta(M+1-z),\\
 R(z,w)&=\frac{\Gamma(1-z-w)\Gamma(1+z)}{\Gamma(1-w)},\\
 V_M(z,w)&=R(z,w)W_M(z+w)
 =\frac{\Gamma(M+1-z-w)\Gamma(1+z)}{\Gamma(1-w)}
                  \zeta(M+1-z-w).
 \label{xid:cld:VW}
\end{align}
In particular $V_M(0,w)=W_M(w)$.  All these functions are holomorphic
in a sufficiently small polydisc about $(0,0)$.

\begin{theorem}[Colliding-grid master theorem]
\label{xid:cld:master}
For $N=r+k\ge1$, the exponential generating function of the
ambient-coordinate finite parts is
\begin{align}
 \mathscr J_{r,k}^{p,q}(z,w)
 &: =\sum_{m,n\ge0}J_{mn;r,k}^{p,q}\frac{z^mw^n}{m!n!}\notag\\
 &=Q^rP^k e^{-Bz-Aw}\left\{
 (-1)^k\frac{V_N(z,w)-e^{Kz}p_r(z)W_N(w)}{z}\right.\notag\\
 &\hspace{35mm}\left.
 +(-1)^r\frac{V_N(w,z)-e^{Kw}p_k(w)W_N(z)}{w}
                                      \right\}.
 \label{xid:cld:master-equation}
\end{align}
Both displayed divided differences have removable denominators separately.
Consequently \eqref{xid:cld:master-equation} is a finite coefficient algorithm
for every $m,n,r,k,p,q$.
\end{theorem}

\begin{proof}
First work where all derivatives of the Hurwitz factors have absolutely
convergent Fourier series.  Orthogonality leaves the frequency indices
$Q\ell$ and $-P\ell$.  The classical coincident Hurwitz product formula
\cite[Theorem 3.1]{EspinosaMoll1}, or the same Fourier calculation followed
by the Riemann functional equation, gives
\begin{align}
 \mathscr C_{r,k}^{p,q}(z,w)
 &:=\int_0^1
    \partial_y^r\zeta(1-z,\{px\})
    \partial_y^k\zeta(1-w,\{qx\})\,dx\notag\\
 &=Q^rP^ke^{-Bz-Aw}
       \left\{(-1)^k\frac{V_N(z,w)}z
                   +(-1)^r\frac{V_N(w,z)}w\right\}.
 \label{xid:cld:spectral}
\end{align}
The integral notation on the left subsequently means meromorphic
continuation.  One can verify the signs directly from
$\partial_y^r\zeta(s,y)=(-1)^r(s)_r\zeta(s+r,y)$ and
$(-1)^r(s)_r\Gamma(1-s-r)=\Gamma(1-s)$.

The local subtraction above continues the integral and identifies its
Taylor coefficients with fixed-coordinate finite parts.  We detail the
terms that distinguish these coefficients from those of
\eqref{xid:cld:spectral}.  At a first-grid point, the singular generating
factor is
\[
 (-1)^r r!p^{-r-1}p_r(z)e^{z\log p}t^{z-r-1}.
\]
Only the Taylor coefficient of degree $r$ of the other factor creates a
pole at $z=0$.  At a common point the other factor is its regular piece
at argument $1$; at a noncommon point it is its ordinary value.  These
arguments, over the first grid, form $d$ copies of
$1/P,2/P,\ldots,1$.  Since $N\ge1$, their trace is
\[
 \frac1p\sum_{b\text{ in first grid}}
      \partial_y^N\mathscr G(w,\eta_b)
   =(-1)^N P^{N-w}W_N(w),
 \qquad \eta_b=\begin{cases}\{qb\},&\{qb\}>0,\\1,&\{qb\}=0.
                         \end{cases}
\]
This is precisely the Hurwitz multiplication formula differentiated
$N$ times in the argument.

The complete $t^{z-1}$ coefficient, summed over the first grid, is
\[
 c_z(z,w)=(-1)^kQ^rP^k e^{z\log p-Aw}p_r(z)W_N(w).
\]
The meromorphic integral of $c_z t^{z-1}$ on $(0,\ell)$ is
$c_z\ell^z/z$.  The generating function of its coefficientwise
ambient finite parts is instead $c_z(\ell^z-1)/z$.
Thus the correction is the \emph{complete} quantity $-c_z(z,w)/z$,
not merely its residue at $z=0$.  The second grid contributes
\[
 -\frac{c_w(z,w)}w,
 \qquad
 c_w(z,w)=(-1)^rQ^rP^k e^{w\log q-Bz}p_k(w)W_N(z).
\]
At common points the singular--singular term is proportional to
$t^{z+w-N-2}$.  Its integral has denominator $z+w-N-1$, which is
nonzero near the origin.  Its analytic integral already equals the
generating function of its finite parts; there is no additional
diagonal contact term at $(0,0)$.

All remaining subtracted integrals converge locally uniformly on a
small spectral polydisc.  Indeed the Taylor remainders gain enough
powers of $t$, and their spectral derivatives introduce only powers
of $\log t$.  They are holomorphic by dominated differentiation.
Subtracting $c_z/z+c_w/w$ from \eqref{xid:cld:spectral} proves
\eqref{xid:cld:master-equation}, since $\log p+B=\log q+A=K$.
If one of $r,k$ is zero, replacing that spectral factor by
$\mathscr G$ adds its scalar $1/z$ or $1/w$; the integral of the other,
positive-order factor is zero by Fourier orthogonality and meromorphic
continuation.  Hence no extra scalar term occurs when $N\ge1$.
\end{proof}

The proof isolates a useful fact: common points introduce a stronger
power divergence, but the finite discrepancy between spectral and
coordinate regularization is governed by the resonant $t^{-1}$ terms.
The stronger merged term must be subtracted, yet its denominator has no
spectral pole at the parameter values being differentiated.

The comparison with rectangular spectral coefficient extraction is
completely explicit.  Define
\[
 \kappa_{m,r}(L)=m![z^{m+1}]p_r(z)e^{Lz},\qquad
 \omega_{n,N}(L)=n![w^n]e^{Lw}W_N(w),
\]
and let $\operatorname{LC}_{mn}\mathscr C$ denote
$m!n![z^mw^n]$ in the rectangular Laurent expansion of
\eqref{xid:cld:spectral}.  Then
\begin{align}
 J_{mn;r,k}^{p,q}-\operatorname{LC}_{mn}\mathscr C_{r,k}^{p,q}
 =-Q^rP^k\{&(-1)^k\kappa_{m,r}(\log p)\omega_{n,N}(-A)\notag\\
           &+(-1)^r\kappa_{n,k}(\log q)\omega_{m,N}(-B)\}.
 \label{xid:cld:spectral-contact-difference}
\end{align}
This finite formula follows by extracting the two complete local
corrections in the proof.  It specifies the difference between the two
regularizations for every index, not only for the zero-index examples.

\begin{corollary}[Finite closure in ordinary zeta derivatives]
For $N\ge1$, every $J_{mn;r,k}^{p,q}$ is a finite linear combination of
$\zeta^{(j)}(N+1)$, $0\le j\le m+n+1$.  Its coefficients lie in the
$\mathbb Q$-algebra generated by $A,B,K$, the rational generalized
harmonic numbers, and positive integer zeta values.  No unevaluated
multiple sum is required.  Writing $M=m+n$, the coefficient of the
highest derivative $\zeta^{(M+1)}(N+1)$ is
\begin{equation}
 (-1)^{M+1}N!Q^rP^k
       \left(\frac{(-1)^k}{m+1}+\frac{(-1)^r}{n+1}\right).
 \label{xid:cld:top-derivative}
\end{equation}
\end{corollary}
\begin{proof}
The finite products $p_r$ and $(1-z)_N$ have rational coefficients.
Moreover
\begin{equation}
 \log R(z,w)=\sum_{j\ge2}\frac{\zeta(j)}j
       \bigl((z+w)^j+(-1)^jz^j-w^j\bigr).
 \label{xid:cld:logR}
\end{equation}
Only total degree at most $m+n+1$ can enter a coefficient of
\eqref{xid:cld:master-equation}.  Taylor expansion of $\zeta(N+1-z-w)$
therefore uses derivatives of order at most $m+n+1$.  The vanishing
constant numerator makes division by $z$ or $w$ an exact shift of
polynomial coefficients.  The highest derivative can arise only from
the degree-$M+1$ term of $\zeta(N+1-z-w)$; the two binomial
coefficients after division by $z$ and $w$ give
\eqref{xid:cld:top-derivative}.  The contact products involve derivatives of
order at most $n$ or $m$, respectively, and cannot alter this coefficient.
\end{proof}

There is a stronger cancellation for equal Stieltjes indices after the
elementary logarithmic shifts forced by the reduced dilations.  Define
\[
 \gamma_{m,L}^{(r)}(x)=\sum_{j=0}^m\binom mj L^{m-j}
                                      \gamma_j^{(r)}(x).
\]
Its generating function is $e^{Lz}\partial_x^r\mathscr G(z,x)$.

\begin{corollary}[Odd-order diagonal collapse]
\label{xid:cld:diagonal-collapse}
If $r+k=N$ is odd, then for every $m\ge0$,
\begin{align}
 &\operatorname{FP}_x\int_0^1
   \gamma_{m,B}^{(r)}(\{px\})\gamma_{m,A}^{(k)}(\{qx\})\,dx\notag\\
 &\qquad=(-1)^kQ^rP^k
       \bigl(\kappa_{m,k}(K)-\kappa_{m,r}(K)\bigr)\omega_{m,N}(0).
 \label{xid:cld:diagonal-collapse-formula}
\end{align}
Thus only zeta derivatives of order at most $m$ remain, whereas a general
pair with these indices can require order $2m+1$.  In particular,
\begin{equation}
 \operatorname{FP}\int_0^1\gamma_m^{(r)}(x)\gamma_m^{(k)}(x)\,dx=0,
 \qquad r+k\ \text{odd},\quad m\ge\max(r,k).
 \label{xid:cld:vanishing-family}
\end{equation}
For $r=0,k=1$ the unequal-dilation identity becomes, for $m\ge1$,
\begin{align}
 &\operatorname{FP}_x\int_0^1
   \gamma_{m,B}(\{px\})\gamma_{m,A}'(\{qx\})\,dx\notag\\
 &\qquad=P(-K)^m\bigl\{\zeta^{(m)}(2)+m\zeta^{(m-1)}(2)\bigr\}.
 \label{xid:cld:diagonal-first-derivative}
\end{align}
Its $m=0$ value is $P\zeta(2)$.
\end{corollary}
\begin{proof}
Multiplying \eqref{xid:cld:master-equation} by $e^{Bz+Aw}$ removes its
external exponential.  Since $(-1)^r=-(-1)^k$, the two spectral terms
$V_N(z,w)/z$ and $V_N(w,z)/w$ have opposite diagonal coefficients.
Only the two explicit contact products remain, giving
\eqref{xid:cld:diagonal-collapse-formula}.  For $p=q=1$, one has $K=0$;
$\kappa_{m,r}(0)=0$ for $m\ge r$ because $p_r$ has degree $r$.
This proves \eqref{xid:cld:vanishing-family}.  Finally $p_1(z)-p_0(z)=-z$
gives $\kappa_{m,1}(K)-\kappa_{m,0}(K)=-K^m$, and
\[
 \omega_{m,1}(0)=(-1)^m
      \{\zeta^{(m)}(2)+m\zeta^{(m-1)}(2)\},\qquad m\ge1.
\]
These two identities prove the last formula.
\end{proof}

\subsection{Polygamma identities and a regularization diagnostic}

Since $\gamma_0=-\psi$, the zero-index specialization is particularly
short.  Put $H_0=0$ and $H_j=\sum_{h=1}^j1/h$.

\begin{corollary}[Every coincident polygamma product]
For $N=r+k\ge1$,
\begin{align}
 &\operatorname{FP}_x\int_0^1
      \psi^{(r)}(\{px\})\psi^{(k)}(\{qx\})\,dx\notag\\
 &\quad=N!Q^rP^k\left\{
 (-1)^{k+1}\bigl[\zeta'(N+1)+(H_N-H_r+K)\zeta(N+1)\bigr]\right.\notag\\
 &\hspace{37mm}\left.
 +(-1)^{r+1}\bigl[\zeta'(N+1)+(H_N-H_k+K)\zeta(N+1)\bigr]
                         \right\}.
 \label{xid:cld:polygamma}
\end{align}
In particular, for odd $N$,
\begin{equation}
 \operatorname{FP}_x\int_0^1
      \psi^{(r)}(\{px\})\psi^{(k)}(\{qx\})\,dx
 =(-1)^{k+1}N!Q^rP^k(H_k-H_r)\zeta(N+1).
 \label{xid:cld:odd-parity}
\end{equation}
These odd-order values are rational multiples of $\pi^{N+1}$ and are
unchanged by a common integer dilation of $p,q$.
\end{corollary}
\begin{proof}
At the origin $W_N(0)=N!\zeta(N+1)$ and
\[
 \partial_z V_N(0,0)
 =-N!\{\zeta'(N+1)+H_N\zeta(N+1)\},
 \qquad p_r'(0)=-H_r.
\]
Take the two removable limits in \eqref{xid:cld:master-equation}.
If $N$ is odd, $r,k$ have opposite parity, so the spectral derivatives
and all scale terms cancel.  Euler's formula for even zeta values
gives the stated rational multiples of $\pi^{N+1}$.
\end{proof}

For example,
\begin{align}
 \operatorname{FP}_x\int_0^1\psi(\{2x\})\psi'(\{3x\})\,dx
      &=\frac{\pi^2}{3},\\
 \operatorname{FP}_x\int_0^1\psi'(\{2x\})\psi'(\{3x\})\,dx
      &=12\{2\zeta'(3)+(1+2\log6)\zeta(3)\},\\
 \operatorname{FP}_x\int_0^1\psi''(\{2x\})\psi'(\{3x\})\,dx
      &=-54\zeta(4)=-\frac{3\pi^4}{5}.
 \label{xid:cld:examples}
\end{align}

The spectral continuation alone gives a different prescription.  Already
for $p=q=1,r=0,k=1$, its rectangular Laurent constant is zero, while
\[
 \operatorname{FP}\int_0^1\psi(x)\psi'(x)\,dx=\zeta(2).
\]
This can be checked without special-function product theory:
$\int_\epsilon^1\psi\psi'=\tfrac12(\psi(1)^2-\psi(\epsilon)^2)$,
and $\psi(\epsilon)=-1/\epsilon-\gamma+\zeta(2)\epsilon+O(\epsilon^2)$.
The constant is $\zeta(2)$.
More generally, substituting $z=\alpha t,w=t$ into that bare spectral
kernel gives constant
$-(\zeta(2)+\zeta'(2))(\alpha-\alpha^{-1})$.
Thus an unspecified one-parameter spectral finite part is not the
ambient-coordinate operation.  The local corrections in
\eqref{xid:cld:master-equation} remove this ambiguity before any coefficient
is extracted.

Two first-index identities illustrate the Stieltjes extension.  With
the same $P,A,K$ as above,
\begin{align}
 J_{10;0,1}^{p,q}
 &=P\left\{\frac{\zeta''(2)}2+K\zeta'(2)
            +\left(\frac{K^2}2+\log p-1\right)\zeta(2)\right\},\\
 J_{01;0,1}^{p,q}
 &=P\left\{-\frac{\zeta''(2)}2-(K+1)\zeta'(2)
            -\left(\frac{K^2}2+A\right)\zeta(2)\right\}.
 \label{xid:cld:first-jets}
\end{align}
Their sum for $p=q=1$ is $-\zeta'(2)-\zeta(2)$, in agreement with
the direct endpoint finite part of $(\gamma_1\gamma_0)'$.

\subsection{The zero-derivative Stieltjes family}

For $r=k=0$, the positive-order kernel must be supplemented by scalar
spectral-pole terms.  It is clearer to give a holomorphic formula using
the undilated coincident kernel, made explicit here:
\[
 \mathscr J_0(z,w)=\sum_{m,n\ge0}J_{mn;0,0}^{1,1}
                         \frac{z^mw^n}{m!n!},
 \qquad g(z)=\zeta(1-z)+\frac1z
          =\sum_{m\ge0}\gamma_m\frac{z^m}{m!}.
\]
The classical coincident spectral kernel is
\begin{equation}
 \mathscr C_{0,0}^{1,1}(z,w)
 =2\Gamma(z)\Gamma(w)(2\pi)^{-z-w}
       \cos\!\left(\frac{\pi(z-w)}2\right)\zeta(z+w).
 \label{xid:cld:N0-spectral}
\end{equation}
The same direct endpoint subtraction gives
\begin{equation}
 \mathscr J_0(z,w)=\mathscr C_{0,0}^{1,1}(z,w)+\frac1{zw}
                             -\frac{g(w)}z-\frac{g(z)}w.
 \label{xid:cld:J0-explicit}
\end{equation}
This grouped expression is holomorphic, as one can also see entirely
algebraically.  Put
\[
 T(z,w)=\frac{wR(z,w)+zR(w,z)}{z+w}.
\]
The numerator vanishes at $w=-z$, because $R(z,-z)=R(-z,z)=1$.
Thus $T$ is holomorphic, and $T(z,0)=T(0,w)=1$.  Applying the Riemann
functional equation to \eqref{xid:cld:N0-spectral} gives
\begin{equation}
 \mathscr J_0(z,w)=
 \frac{1-T(z,w)+(z+w)T(z,w)g(z+w)-wg(w)-zg(z)}{zw}.
 \label{xid:cld:J0-removable}
\end{equation}
The numerator vanishes on both coordinate axes, so division by $zw$
is removable.  This provides a self-contained finite coefficient rule
for the undilated zero-derivative case, without a previously tabulated
coefficient array.

Write $L_p=\log p$ and $L_q=\log q$.  Direct local subtraction, exactly
as in the master theorem, gives
\begin{align}
 \mathscr J_{0,0}^{p,q}(z,w)
 ={}&e^{-Bz-Aw}\mathscr J_0(z,w)\notag\\
 &+e^{-Aw}\frac{e^{-Bz}-e^{L_pz}}z\,g(w)
  +e^{-Bz}\frac{e^{-Aw}-e^{L_qw}}w\,g(z)\notag\\
 &+\frac{1-e^{-Bz-Aw}-e^{L_pz}+e^{L_pz-Aw}
                       -e^{L_qw}+e^{L_qw-Bz}}{zw}.
 \label{xid:cld:N0-kernel}
\end{align}
Every denominator on the right is removable; the last numerator vanishes
on both coordinate axes.  Equivalently, the full separated-grid
coefficient formula of the incoming report remains valid at collision
after replacing $I_{ij}(\theta)$ by $J_{ij;0,0}^{1,1}$ and both single
Stieltjes arguments $\theta,1-\theta$ by $1$.  This substitution is a
conclusion of local subtraction, not evaluation of a divergent integrand
at a singular shift.

For completeness, the scalar correction behind
\eqref{xid:cld:N0-kernel} is visible before extracting coefficients:
\[
 \mathscr J_{0,0}^{p,q}
 =\mathscr C_{0,0}^{p,q}+\frac1{zw}
  -\frac{e^{L_pz}}z\left(e^{-Aw}\zeta(1-w)+\frac1w\right)
  -\frac{e^{L_qw}}w\left(e^{-Bz}\zeta(1-z)+\frac1z\right).
\]
The terms in parentheses are holomorphic by cancellation of the zeta
pole.  The additional $1/(zw)$ comes from the two scalar terms used to
define $\mathscr G$.  Expanding this identity and its undilated version
proves \eqref{xid:cld:N0-kernel}.

In particular,
\begin{equation}
 \operatorname{FP}_x\int_0^1\psi(\{px\})\psi(\{qx\})\,dx
 =2\gamma_1-2\zeta(2)-2\gamma K-K^2+\log p\log q.
 \label{xid:cld:digamma}
\end{equation}
Indeed, \eqref{xid:cld:logR} gives
$T(z,w)=1+2\zeta(2)zw+O((|z|+|w|)^3)$, so
\eqref{xid:cld:J0-removable} yields $J_{00;0,0}^{1,1}=2\gamma_1-2\zeta(2)$.
This recovers the undilated coincident identity in the incoming report.
The occurrence of $K$ in the colliding unequal case is forced by the
ambient subtraction scale.

\subsection{Exact matching with the noncollision expansion}

Let
\[
 F_{mn;r,k}^{p,q}(\theta)=\FP_x\int_0^1
 \gamma_m^{(r)}(\{px\})\gamma_n^{(k)}(\{qx+a\})\,dx,
 \qquad \theta=\{Pa\}\in(0,1),
\]
denote the separated finite part. Its dependence on $a$ factors through
$\theta$ by the finite covering formula, or by Fourier orthogonality.  The next theorem resolves the collision
matching question, including every mixed derivative and Stieltjes index.
Set
\[
 C_N(z,w)=\frac{\Gamma(N+1-z-w)\Gamma(1+z)}{\Gamma(1-w)}.
\]
Unlike $V_N$, this definition is also regular for $N=0$.

\begin{theorem}[Convergent collision expansion]
\label{xid:cld:collision-expansion}
For $m,n,r,k\ge0$, define the explicit logarithmic polynomial
\begin{align}
 \mathcal P_{mn;r,k}^{p,q}(L)
 ={}&(-1)^kQ^rP^k m!n![z^mw^n]\notag\\
 &\quad e^{-Bz-Aw+wL}
       \frac{C_N(z,w)e^{zL}-e^{Kz}p_r(z)(1-w)_N}{z}.
 \label{xid:cld:collision-polynomial}
\end{align}
The quotient is holomorphic.  The polynomial has degree $m+n+1$ and
leading coefficient
\[
 \frac{(-1)^kN!Q^rP^k}{m+1}.
\]
There is an exact expansion
\begin{equation}
 F_{mn;r,k}^{p,q}(\theta)
 =\frac{\mathcal P_{mn;r,k}^{p,q}(\log\theta)}{\theta^{N+1}}
       +J_{mn;r,k}^{p,q}+\sum_{j\ge1}h_j\theta^j,
 \qquad 0<\theta<1,
 \label{xid:cld:collision-expansion-equation}
\end{equation}
whose regular power series converges locally uniformly for $|\theta|<1$.
Thus the regular constant at collision agrees exactly with the
independently defined merged-grid finite part.
\end{theorem}

\begin{proof}
The same Fourier and local-subtraction proof gives the separated kernel
by replacing $V_N(z,w)$ with
$C_N(z,w)\zeta(N+1-z-w,\theta)$ and $W_N(w)$ with
$(1-w)_N\zeta(N+1-w,\theta)$ in the first divided difference, and
by making the analogous replacements with $1-\theta$ in the second.
When $N=0$, the additional pure scalar terms containing $1/(zw)$
are independent of $\theta$ and remain unchanged; the Hurwitz terms
inside the contact numerators take the same $\theta$ or $1-\theta$
replacements. These statements follow directly from the scalar formula
preceding \eqref{xid:cld:digamma}.

Use the exact shift and Taylor identities
\begin{align*}
 \zeta(s,\theta)&=\theta^{-s}+\zeta(s,1+\theta),\\
 \zeta(s,1+\theta)&=\sum_{j\ge0}\frac{(-1)^j(s)_j}{j!}
                                      \zeta(s+j)\theta^j,\\
 \zeta(s,1-\theta)&=\sum_{j\ge0}\frac{(s)_j}{j!}
                                      \zeta(s+j)\theta^j.
\end{align*}
The first isolated power produces exactly
\eqref{xid:cld:collision-polynomial}.  The $j=0$ regular terms give the
colliding kernel proved above.  For $j\ge1$ there is no zeta pole near
the spectral origin, and the coefficient-generating function is
\begin{align}
 \mathscr H_j(z,w)
 ={}&\frac{Q^rP^ke^{-Bz-Aw}}{j!}
   \left\{(-1)^{k+j}
      \frac{V_{N+j}(z,w)-e^{Kz}p_r(z)W_{N+j}(w)}z\right.\notag\\
 &\hspace{30mm}\left.
    +(-1)^r
      \frac{V_{N+j}(w,z)-e^{Kw}p_k(w)W_{N+j}(z)}w\right\}.
 \label{xid:cld:regular-coefficients}
\end{align}
Here $h_j=m!n![z^mw^n]\mathscr H_j(z,w)$.
After division by $j!$, the gamma factors and rising factorials grow
at most polynomially in $j$ on a fixed sufficiently small spectral polydisc; their spectral
derivatives introduce powers of $\log j$.  Thus $|\theta|^j$ gives
locally uniform convergence for $|\theta|<1$, before and after any fixed
coefficient extraction.  For $N=0$, the common $j=0$ spectral pole is
removed in the complete kernel \eqref{xid:cld:N0-kernel}; it does not affect
the convergent $j\ge1$ series.

Finally the highest $L$-degree in
\eqref{xid:cld:collision-polynomial} comes from
$C_N(0,0)=N!$ and $e^{(z+w)L}/z$.  The second term in that numerator
contains $L$ only through $e^{wL}$ and cannot contribute degree
$m+n+1$.  Coefficient extraction gives the asserted leading coefficient.
\end{proof}

The lowest polynomial is already informative:
\begin{equation}
 \mathcal P_{00;r,k}^{p,q}(L)
    =(-1)^kN!Q^rP^k(L-H_N+H_r-K).
 \label{xid:cld:lowest-collision}
\end{equation}
For example,
\[
 \lim_{\theta\downarrow0}\left\{
 \operatorname{FP}_x\int_0^1
   \psi(\{2x\})\psi'(\{3x+a\})\,dx
     +\frac{2(\log\theta-1-\log6)}{\theta^2}\right\}
 =\frac{\pi^2}{3},\qquad \theta=\{2a\}.
\]
The finite term is therefore fixed by an exact convergent expansion,
including the ambient scale and every local subtraction term.

\subsection{Common dilation and an integration safeguard}

For $N\ge1$ and a further positive integer $e$, put $D=\log e$.  The reduced pair
$P,Q$ is unchanged and $K$ becomes $K+D$.  From the master kernel,
\begin{align}
 \mathscr J_{r,k}^{ep,eq}-\mathscr J_{r,k}^{p,q}
 ={}&-Q^rP^ke^{-Bz-Aw}
   \left\{(-1)^k p_r(z)W_N(w)e^{Kz}\frac{e^{Dz}-1}{z}\right.\notag\\
 &\hspace{30mm}\left.
        +(-1)^r p_k(w)W_N(z)e^{Kw}\frac{e^{Dw}-1}{w}\right\}.
 \label{xid:cld:common-scale}
\end{align}
This is a finite polynomial in $D$ at every fixed pair of Stieltjes
indices.  At zero index and odd $N$, the two terms cancel, explaining
the exact scale independence in \eqref{xid:cld:odd-parity}.  At even $N$ the
change is the negative of $D$ times the summed local $t^{-1}$ coefficient.

The incoming separated-grid theorem is valid as stated and requires no
retraction.  Its collision question can now be marked resolved by
\eqref{xid:cld:master-equation}, \eqref{xid:cld:N0-kernel}, and
\eqref{xid:cld:collision-expansion-equation}.  What must not be done during
integration is to multiply the two grid-supported contact distributions
at collision, or to keep only residues in the local spectral corrections.
The pointwise product has the unambiguous ambient finite part proved
here; its analytic representation requires the complete pole numerators.

% Verification: verify_collision.py evaluates the local-subtraction integral
% independently for nine polygamma cases; collision_validation.json records
% the floating-point residuals.  These are consistency checks, not enclosures.

% END sections/collision.tex


% BEGIN sections/nonlinear.tex
\section{Nonlinear endpoint coordinates: a finite exact law}
\label{xid:sec:nonlinear}

The two recent dilation reports explicitly ask for the replacement of
their affine coordinate $pt$ by $pt+q_2t^2+q_3t^3+\cdots$
\cite{IncomingMellin,IncomingIdentities}.  This section answers that
question for every Stieltjes index and every argument derivative.  The
answer is local: a singularity at any other point is treated after
translation.  All finite parts use a right-hand cutoff in the indicated
coordinate.  Two-sided symmetric prescriptions are different objects.

Let $f$ be a real analytic local diffeomorphism with
\begin{equation}\label{xid:eq:f-germ}
 f(0)=0,\qquad f'(0)=p>0,\qquad
 H_f(t)=\frac{f(t)}t=p+q_2t+q_3t^2+\cdots.
\end{equation}
Choose a sufficiently small neighborhood so that $H_f$ is positive on
the real axis.  The power $H_f(t)^z$ uses the analytic logarithm whose
value at zero is $\log p$.  For a distribution $T$ the scalar pullback is
\begin{equation}\label{xid:eq:pullback-definition}
 \langle f^*T,\varphi\rangle
 =\left\langle T,
 \frac{\varphi(f^{-1}(y))}{f'(f^{-1}(y))}\right\rangle.
\end{equation}
Thus $f^*T=T\circ f$ for ordinary functions.  The Jacobian in this
definition is essential.

\subsection{A general coefficient lemma}

Write $\delta^{(j)}$ for the $j$th distributional derivative at zero,
so that $\langle\delta^{(j)},\varphi\rangle=(-1)^j\varphi^{(j)}(0)$.
For integers $q\ge1$ and $n\ge0$ set
\[
 U_{q,n}=\FP_y\bigl[\mathbf1_{y>0}y^{-q}\log^n y\bigr].
\]
Only the germ at zero matters; a smooth cutoff may be inserted away
from zero in every formula below.

\begin{theorem}[Coordinate defect for a logarithmic pole]
\label{xid:thm:coordinate-monomial}
One has
\begin{equation}\label{xid:eq:coordinate-monomial}
 f^*U_{q,n}
 =\FP_t\bigl[\mathbf1_{t>0}f(t)^{-q}\log^n f(t)\bigr]
       +\mathcal C_f(q,n),
\end{equation}
where the complete correction is the finite distribution
\begin{equation}\label{xid:eq:C-generator}
 \mathcal C_f(q,n)
 =n!\sum_{j=0}^{q-1}\frac{(-1)^j}{j!}
   [z^{n+1}t^{q-1-j}]H_f(t)^{z-q}\,\delta^{(j)}.
\end{equation}
It depends only on $f'(0),\ldots,f^{(q)}(0)$.
\end{theorem}

\begin{proof}
Begin with the locally integrable meromorphic family
$y_+^{z-q}$ for $\Re z>q-1$.  Pairing its pullback with a test function
gives
\[
 \int_0^h t^{z-q}H_f(t)^{z-q}\varphi(t)\,dt.
\]
Subtract the Taylor polynomial in $t$ through degree $q-1$ from the
holomorphic factor $H_f(t)^{z-q}\varphi(t)$.  The remainder is locally
normally integrable for $\Re z>-1$.  A Taylor term of degree $a$ has
integral
\[
 g_a(z)\frac{h^{z-q+a+1}}{z-q+a+1}.
\]
Near $z=0$, only $a=q-1$ has a pole.  The other Taylor terms have
nonzero denominators at zero, and taking their spectral coefficients
commutes with the coordinate finite part.  For the exceptional term,
the cutoff expression is
\[
 g_{q-1}(z)\frac{h^z-\varepsilon^z}{z}.
\]
The constant term in $\varepsilon$ of its coefficient of $z^n$ is the
regular coefficient of $g_{q-1}(z)h^z/z$, less
$[z^{n+1}]g_{q-1}(z)$.  In the original $y$ coordinate that last
correction is zero, because its Taylor coefficient is independent of
$z$.  Therefore the difference after pullback is
\[
 n![z^{n+1}t^{q-1}]H_f(t)^{z-q}\varphi(t).
\]
Expanding the test function and using the convention for
$\delta^{(j)}$ proves \eqref{xid:eq:C-generator}.  Its largest required
power of $t$ is $q-1$, proving the finite dependence assertion.
\end{proof}

Every function with a finite expansion in negative integer powers and
nonnegative logarithmic powers is covered by linearity.  Terms
$t^a\log^b t$ with $a\ge0$ are locally integrable and create no defect.
Consequently the theorem also applies to products of singular germs
after their ordinary local expansion has been performed; it does not
require a multiplication of distributions.

\begin{corollary}[A sharp disappearance threshold]
\label{xid:cor:monomial-threshold}
If $f'(0)=1$, then $\mathcal C_f(q,n)=0$ for $n\ge q-1$.
For $q\ge2$, writing $f(t)=t+a t^2+O(t^3)$, one has
\begin{equation}\label{xid:eq:monomial-top}
 \mathcal C_f(q,q-2)=\frac{a^{q-1}}{q-1}\,\delta.
\end{equation}
Thus the threshold is attained when $a\ne0$.
\end{corollary}

\begin{proof}
When $H_f(0)=1$, $[t^k]H_f(t)^{z-q}$ is a polynomial in $z$ of degree
at most $k$.  This follows by expanding
$\exp((z-q)\log H_f(t))$, since $\log H_f(t)=O(t)$.
In \eqref{xid:eq:C-generator}, $k\le q-1$, so its coefficient of
$z^{n+1}$ vanishes for $n\ge q-1$.  At $n=q-2$, only $j=0$ can
survive; the leading coefficient of $[t^{q-1}]H_f(t)^{z-q}$ is
$a^{q-1}/(q-1)!$.  Multiplication by $(q-2)!$ proves the formula.
\end{proof}

\subsection{Composition is an exact cocycle}

For any logarithmic singular germ $u$ define
\[
 C_f(u)=f^*(\FP u)-\FP(u\circ f).
\]
Theorem~\ref{xid:thm:coordinate-monomial} computes this distribution from the
finite singular expansion of $u$.  The expansion of $u\circ f$ is
again of the same kind, up to locally integrable terms.

\begin{theorem}[Composition and explicit delta pullback]
\label{xid:thm:cocycle}
For orientation-preserving analytic local coordinates $f,g$,
\begin{equation}\label{xid:eq:cocycle}
 C_{f\circ g}(u)=g^*C_f(u)+C_g(u\circ f).
\end{equation}
The delta derivatives in this formula can be computed without series
reversion:
\begin{equation}\label{xid:eq:delta-pullback}
 f^*\delta^{(r)}
 =\sum_{j=0}^r\frac{(-1)^{r-j}r!}{j!}
     [t^{r-j}]H_f(t)^{-r-1}\,\delta^{(j)}.
\end{equation}
\end{theorem}

\begin{proof}
Pullbacks satisfy $(f\circ g)^*=g^*f^*$.  Insert and subtract
$g^*\FP(u\circ f)$ in the definition of $C_{f\circ g}$; this gives
\eqref{xid:eq:cocycle} exactly.  For \eqref{xid:eq:delta-pullback}, take the
residue at $z=0$ in the proof of Theorem~\ref{xid:thm:coordinate-monomial}
with $q=r+1$.  The residue of $y_+^{z-r-1}$ is
$(-1)^r\delta^{(r)}/r!$.  The residue after pullback pairs with
$\varphi$ as $[t^r]H_f(t)^{-r-1}\varphi(t)$.
Expanding $\varphi$ gives the displayed formula.
\end{proof}

This proves the composition requirement in the source question, rather
than merely checking it for low-degree maps.  The accompanying exact
code also compares both sides using independently composed rational
Taylor series.

\subsection{Every Stieltjes index and every argument derivative}

Let $G_n$ be the canonical periodic finite part of $\gamma_n$ in the
unit coordinate.  Its local meromorphic generator is
\begin{equation}\label{xid:eq:periodic-G-local}
 Z(1-z)=\frac{\delta_0-1}{z}
          +\sum_{n\ge0}G_n\frac{z^n}{n!}.
\end{equation}
Indeed, locally across the endpoint,
$\zeta(1-z,\{y\})=y_+^{z-1}+\zeta(1-z,1+y)$.
The second term is smooth in $y$ and has only the constant spectral
pole $-1/z$.  Set
\begin{equation}\label{xid:eq:Er}
 E_r(z)=\prod_{h=1}^r(1-z/h),\qquad E_0(z)=1.
\end{equation}

\begin{theorem}[Nonlinear Stieltjes contact law]
\label{xid:thm:stieltjes-nonlinear}
For all $n,r\ge0$, as local distributions,
\begin{equation}\label{xid:eq:stieltjes-nonlinear}
 f^*(D_y^rG_n)
 =\FP_t\bigl[\gamma_n^{(r)}(\{f(t)\})\bigr]
       +A_{n,r}(f),
\end{equation}
where
\begin{equation}\label{xid:eq:A-generator}
 A_{n,r}(f)=n!r!\sum_{j=0}^r\frac{(-1)^{r+j}}{j!}
 [z^{n+1}t^{r-j}]
 E_r(z)H_f(t)^{z-r-1}\,\delta^{(j)}.
\end{equation}
The correction depends only on the $(r+1)$-jet of $f$.
\end{theorem}

\begin{proof}
The singular part of $D_y^r Z(1-z)$ is
$(-1)^r(1-z)_r y_+^{z-r-1}$.  After pullback, pair it with a test
function and repeat the Taylor subtraction proof above.  Only its
degree-$r$ Taylor coefficient creates a pole at $z=0$; the correction
is
\[
 n![z^{n+1}t^r](-1)^r(1-z)_r
                  H_f(t)^{z-r-1}\varphi(t).
\]
Since $(1-z)_r=r!E_r(z)$, this is \eqref{xid:eq:A-generator}.
The smooth local Hurwitz term commutes with coefficient extraction
and contributes no contact distribution.  The largest Taylor order
needed is $r$, which proves the jet assertion.
\end{proof}

For an affine germ $f(t)=pt$, only the term $j=r$ remains, and
\begin{equation}\label{xid:eq:A-affine}
 A_{n,r}(pt)=p^{-r-1}b_{n,r}(\log p)\delta^{(r)},\qquad
 b_{n,r}(L)=n![z^{n+1}]e^{Lz}E_r(z).
\end{equation}
After multiplication by $p^r$, this is precisely the earlier
ambient-derivative dilation law.  In particular,
\[
 b_{0,r}(L)=L-H_r,\qquad
 b_{1,r}(L)=\frac{L^2}{2}-H_rL+
                    \frac{H_r^2-H_r^{(2)}}2.
\]
The coefficient algebra is therefore an explicit extension of the
source's harmonic contact polynomial, with the same normalization.

For $f(t)=pt+q_2t^2+O(t^3)$ and $L=\log p$, the first two argument
orders take the particularly small form
\begin{align}
 A_{n,0}(f)&=\frac{L^{n+1}}{(n+1)p}\delta,\label{xid:eq:A0}\\
 A_{n,1}(f)&=\frac{n!}{p^2}[z^{n+1}](1-z)e^{Lz}\delta'
 +\frac{n!q_2}{p^3}[z^{n+1}](2-3z+z^2)e^{Lz}\delta.
 \label{xid:eq:A1}
\end{align}
Curvature creates an actual lower-order contact term; replacing $f$
by its tangent line loses it.

\begin{corollary}[Sharp unit-tangent Stieltjes threshold]
\label{xid:cor:stieltjes-threshold}
Suppose $f(t)=t+a t^2+O(t^3)$ and $r\ge1$.  Then
\begin{equation}\label{xid:eq:stieltjes-threshold}
 A_{n,r}(f)=0\quad(n\ge2r),\qquad
 A_{2r-1,r}(f)=\frac{(2r-1)!}{r!}a^r\delta.
\end{equation}
\end{corollary}

\begin{proof}
The degree in $z$ of
$[t^{r-j}]E_r(z)H_f(t)^{z-r-1}$ is at most $2r-j$.
This proves the vanishing assertion.  At $n=2r-1$ only $j=0$
survives.  The leading coefficients are $(-1)^r/r!$ for $E_r$
and $a^r/r!$ for $[t^r]H_f^{z-r-1}$.  Substitute them into
\eqref{xid:eq:A-generator} to obtain the second formula.
\end{proof}

For $f(t)=t+a t^2+b t^3+O(t^4)$ some exact coefficients are
\begin{align}
 A_{0,1}&=-\delta'-3a\delta,& A_{1,1}&=a\delta,\notag\\
 A_{0,2}&=-\tfrac32\delta''-11a\delta'
                       +(11b-25a^2)\delta,&
 A_{3,2}&=3a^2\delta.\label{xid:eq:A-table}
\end{align}
Together with \eqref{xid:eq:stieltjes-threshold}, these examples show why
the earlier zero contact terms at unit affine scale cannot simply be
transported through a nonlinear coordinate.

\subsection{Exact antiderivatives and their endpoint constants}

There is a useful scalar consequence which does not require any
distribution notation.  Suppose $f$ is increasing on $[0,b]$,
$f(0)=0$, $f'(0)>0$, and $B=f(b)>0$.  Assume that the Taylor germ at
zero is as above.  For $r\ge1$ define
\begin{equation}\label{xid:eq:primitive-E}
 \mathcal E_{n,r}(f)=\gamma_n^{(r-1)}(1)
 +(-1)^{r-1}n![z^nt^r](1-z)_{r-1}H_f(t)^{z-r}.
\end{equation}

\begin{theorem}[Finite-part primitive endpoint identity]
\label{xid:thm:nonlinear-primitive}
For every $n\ge0$ and $r\ge1$,
\begin{equation}\label{xid:eq:nonlinear-primitive}
 \FP_t\int_0^b f'(t)\gamma_n^{(r)}(f(t))\,dt
       =\gamma_n^{(r-1)}(B)-\mathcal E_{n,r}(f).
\end{equation}
For argument order zero the exact change-of-coordinate identity is
\begin{equation}\label{xid:eq:primitive-order0}
 \FP_t\int_0^b f'(t)\gamma_n(f(t))\,dt
 =\FP_y\int_0^B\gamma_n(y)\,dy
       -\frac{\log^{n+1}p}{n+1}.
\end{equation}
\end{theorem}

\begin{proof}
At a positive cutoff, the fundamental theorem of calculus gives
$\gamma_n^{(r-1)}(B)-\gamma_n^{(r-1)}(f(\varepsilon))$.
Use
\[
 \gamma_n(y)=y^{-1}\log^n y+\gamma_n(1+y)
\]
and differentiate $r-1$ times.  The smooth term has constant
$\gamma_n^{(r-1)}(1)$.  Write its singular term as the coefficient
of $z^n/n!$ in
$(-1)^{r-1}(1-z)_{r-1}t^{z-r}H_f(t)^{z-r}$.
The constant term in $t$ comes only from the Taylor term of degree
$r$, with no power of $\log t$ retained.  This is exactly
\eqref{xid:eq:primitive-E}.  For $r=0$, substitution gives a lower
endpoint primitive $\log^{n+1}f(\varepsilon)/(n+1)$, whose constant
term is $\log^{n+1}p/(n+1)$.  All smooth endpoint contributions agree
in the two coordinates, proving \eqref{xid:eq:primitive-order0}.
\end{proof}

For example, put $f(t)=t+c t^2$ on an interval where it is increasing.
Then $B=b+cb^2$ and the formulas give
\begin{align}
 \FP\int_0^b(1+2ct)\psi'(t+ct^2)\,dt
    &=\psi(B)+\gamma-c,\label{xid:eq:psi-prime-curvature}\\
 \FP\int_0^b(1+2ct)\psi''(t+ct^2)\,dt
    &=\psi'(B)-\zeta(2)-3c^2,\label{xid:eq:psi-double-curvature}\\
 \FP\int_0^b(1+2ct)\gamma_1'(t+ct^2)\,dt
    &=\gamma_1(B)-\gamma_1(1)-c.\label{xid:eq:gamma1-curvature}
\end{align}
These are exact antiderivative identities with explicit curvature
constants.  In particular, a formal substitution which silently
equates the $t$ and $f(t)$ finite parts gives an incorrect constant.

The zero-order Hurwitz primitive on the right of
\eqref{xid:eq:primitive-order0} can be written, if desired, by integrating
$\zeta(1-z,y)$ as $\zeta(-z,y)/z$ and taking its regular spectral
coefficients.  At $n=0$, $\gamma_0=-\psi$ gives the familiar
logarithmic Gamma primitive.  The new ingredient above is the exact
endpoint transformation, not the classical Hurwitz primitive itself
\cite{EspinosaMoll2}.

% END sections/nonlinear.tex


% BEGIN sections/triple_dilation.tex
\section{Three unequal frequencies and all their regularized jets}
\label{xid:sec:triple-dilation}

The current incoming report treats three distinct translates and asks
for three unequal integer frequencies \cite{IncomingIdentities}.
Fourier orthogonality leaves a lattice of rank two in this problem.
Finite congruence filters nevertheless give an explicit finite family
of colored Tornheim kernels.  Local subtraction then converts the
spectral identity into an all-index Stieltjes identity in the ambient
coordinate.  The classical use of colored Tornheim sums and their
polylogarithmic reductions is credited to the literature
\cite{ZhaoColored}; the formulas here address the unequal-frequency
and fixed-coordinate question from the supplied research.

Let $p_1,p_2,p_3$ be positive integers and $a_1,a_2,a_3$ be real
shifts, considered modulo one.  Write
\[
 S_i=\{x\in\mathbb R/\mathbb Z:p_i x+a_i\in\mathbb Z\},
 \qquad e(t)=e^{2\pi i t}.
\]
Throughout this section assume that $S_1,S_2,S_3$ are pairwise
disjoint.  Equivalently, for every $i\ne j$,
\begin{equation}\label{xid:eq:triple-separation}
 p_j a_i-p_i a_j\notin\gcd(p_i,p_j)\mathbb Z.
\end{equation}
At each grid point the cutoff is the right ambient coordinate
$t=x-x_0$; the cutoffs at different grid points can tend to zero
independently.

\subsection{Weighted colored kernels and finite character filters}

For positive integers $a,b$ and nontrivial unit complex numbers $X,Y$
define initially in an absolutely convergent half-space
\begin{equation}\label{xid:eq:weighted-T}
 T_{a,b}(A,B,C;X,Y)
 =\sum_{m,n\ge1}\frac{X^mY^n}{m^A n^B(am+bn)^C}.
\end{equation}
The positive real denominators use real logarithms.  For $\Re C>0$
the entire continuation in $A,B$ is given by
\begin{equation}\label{xid:eq:weighted-T-mellin}
 T_{a,b}(A,B,C;X,Y)
 =\frac1{\Gamma(C)}\int_0^\infty t^{C-1}
        \Li_A(Xe^{-at})\Li_B(Ye^{-bt})\,dt.
\end{equation}
In fact this is jointly entire in all three orders after continuation.
To prove the statement, subtract any finite Taylor polynomial at
$t=0$ from the product of polylogarithms.  It is holomorphic there,
locally uniformly in $A,B$, because $X,Y\ne1$.  The Taylor integrals
have poles only at $C=0,-1,\ldots$, each canceled by $1/\Gamma(C)$.
The integral over $[1,\infty)$ converges locally normally for all
orders because its integrand decays exponentially.  Increasing the
Taylor order proves entire continuation and licenses every fixed
mixed order derivative.

These weighted kernels introduce no extra class of functions beyond
ordinary colored Tornheim kernels.  Choose any roots $\xi^a=X$ and
$\eta^b=Y$.  Applying two finite root filters gives
\begin{equation}\label{xid:eq:weighted-T-reduction}
 T_{a,b}(A,B,C;X,Y)
 =a^{A-1}b^{B-1}
 \sum_{r=0}^{a-1}\sum_{s=0}^{b-1}
 T_{1,1}(A,B,C;\xi e(r/a),\eta e(s/b)).
\end{equation}
Different choices of roots merely permute the summands.  None of
these twists is one, since $X,Y\ne1$.  The proof first changes
variables $u=am,v=bn$ in the absolutely convergent series, and then
uses $a^{-1}\sum_r e(ru/a)=\mathbf1_{a\mid u}$ and its $b$ analogue.
Analytic continuation proves the identity at all orders.

For each $k\in\{1,2,3\}$ let $i<j$ be the remaining indices, and
for $h=0,\ldots,p_k-1$ put
\begin{align}
 X_{ik,h}&=e\left(a_i-\frac{p_i a_k}{p_k}
                              +\frac{hp_i}{p_k}\right),\notag\\
 X_{jk,h}&=e\left(a_j-\frac{p_j a_k}{p_k}
                              +\frac{hp_j}{p_k}\right),\label{xid:eq:triple-twists}\\
 \vartheta_k&=\frac\pi2(\lambda_i+\lambda_j-\lambda_k).
 \notag
\end{align}
Every twist in \eqref{xid:eq:triple-twists} is nontrivial.  Indeed,
$X_{ik,h}=1$ for some integer $h$ is equivalent to
$p_k a_i-p_i a_k\in\gcd(p_i,p_k)\mathbb Z$, which contradicts
\eqref{xid:eq:triple-separation}.

\begin{theorem}[Finite spectral formula for three unequal frequencies]
\label{xid:thm:triple-spectral}
For $\Re\lambda_i>0$,
\begin{equation}\label{xid:eq:triple-K}
 K(\boldsymbol\lambda;\boldsymbol p,\boldsymbol a)
 =\int_0^1\prod_{i=1}^3
       \zeta(1-\lambda_i,\{p_i x+a_i\})\,dx
 =\frac{\prod_i\Gamma(\lambda_i)}{(2\pi)^{\sum_i\lambda_i}}
       \mathcal S(\boldsymbol\lambda),
\end{equation}
where the finite sum is
\begin{align}
 \mathcal S(\boldsymbol\lambda)
 =\sum_{k=1}^3p_k^{\lambda_k-1}\sum_{h=0}^{p_k-1}
 \bigl\{&e^{-i\vartheta_k}
 T_{p_i,p_j}(\lambda_i,\lambda_j,\lambda_k;
                         X_{ik,h},X_{jk,h})\notag\\
 &+e^{i\vartheta_k}
 T_{p_i,p_j}(\lambda_i,\lambda_j,\lambda_k;
                         X_{ik,h}^{-1},X_{jk,h}^{-1})\bigr\}.
 \label{xid:eq:triple-cones}
\end{align}
The right side is a meromorphic continuation to all three order
variables.  It has $2(p_1+p_2+p_3)$ weighted-kernel terms before
symmetry or character simplifications.
\end{theorem}

\begin{proof}
Start with $\Re\lambda_i>1$, where each Hurwitz Fourier series is
absolutely convergent.  For a nonzero frequency $n$, its coefficient
is
\[
 \frac{\Gamma(\lambda_i)}{(2\pi|n|)^{\lambda_i}}
      e^{-i\pi\lambda_i\operatorname{sgn}(n)/2}.
\]
After dilation and translation, integration retains precisely
$p_1n_1+p_2n_2+p_3n_3=0$, with all $n_i\ne0$.
There are six disjoint sign cones.  In the cone with $n_i=m>0$,
$n_j=n>0$ and $n_k<0$, one has
\[
 -n_k=\frac{p_i m+p_j n}{p_k}\in\mathbb N.
\]
The denominator contributes $p_k^{\lambda_k}$ and the phase of
the Fourier coefficient is $e^{-i\vartheta_k}$.  Insert
\[
 \mathbf1_{p_k\mid p_i m+p_j n}
   =\frac1{p_k}\sum_{h=0}^{p_k-1}
                  e\left(\frac{h(p_i m+p_j n)}{p_k}\right).
\]
Combining this filter with the translation phases produces exactly
the first term in braces in \eqref{xid:eq:triple-cones}, including the
factor $p_k^{\lambda_k-1}$.  Reversing all frequency signs gives
the second term, after reindexing the finite filter by $h\mapsto-h$.
Each lattice point occurs once.

The original integral is jointly holomorphic for $\Re\lambda_i>0$:
near a grid point only one factor is singular, with a locally
integrable majorant $t^{\sigma-1}$ on any compact parameter set.
Order derivatives add only logarithmic powers.  Entire continuation
of the colored kernels makes the right side holomorphic in this
half-space as well.  The identity theorem proves the asserted
extension, and the Gamma factors give the meromorphic continuation.
\end{proof}

There is also a longer formula directly in the known equal-frequency
kernel.  The ordinary Hurwitz multiplication law gives
\begin{equation}\label{xid:eq:covering-triple}
 K(\boldsymbol\lambda;\boldsymbol p,\boldsymbol a)
 =\prod_i p_i^{\lambda_i-1}
 \sum_{h_1=0}^{p_1-1}\sum_{h_2=0}^{p_2-1}\sum_{h_3=0}^{p_3-1}
 K\left(\boldsymbol\lambda;(1,1,1),
        \left(\frac{a_i+h_i}{p_i}\right)_{i=1}^3\right).
\end{equation}
All three translated endpoints in every summand are distinct.  This
provides an alternative finite reduction and confirms that the rank-two
lattice does not require a new, unbounded family of kernels.

\subsection{All Stieltjes and argument jets in the ambient coordinate}

Set $\mathscr G(\lambda,y)=\zeta(1-\lambda,y)+1/\lambda$ as before,
and define the pair kernels by
\[
 C_{ij}(\lambda_i,\lambda_j)
 =\int_0^1\zeta(1-\lambda_i,\{p_ix+a_i\})
                \zeta(1-\lambda_j,\{p_jx+a_j\})\,dx,
\]
with their ordinary separated-grid analytic continuation.
They have the finite Gamma--polylogarithm expression
\eqref{xid:eq:pair-spectral}. All shift derivatives below are taken within
a local chamber of separated real shifts. Define
\begin{equation}\label{xid:eq:triple-regular-spectral}
 \mathcal L(\boldsymbol\lambda)
 =K(\boldsymbol\lambda)+\sum_{\substack{\{i,j,k\}=\{1,2,3\}\\i<j}}
       \frac{C_{ij}(\lambda_i,\lambda_j)}{\lambda_k}
       +\frac1{\lambda_1\lambda_2\lambda_3}.
\end{equation}
In an ordinary convergent domain this is the integral of the product
of the three $\mathscr G$ functions.  The terms with just one Hurwitz
factor vanish because its integral over a period is zero.

For $\boldsymbol r=(r_1,r_2,r_3)\in\mathbb N_0^3$ let
$\partial_{\boldsymbol a}^{\boldsymbol r}=\prod_i\partial_{a_i}^{r_i}$.
For $x_0\in S_i$ and $j\ne i$ put
$\eta_{ij}(x_0)=\{p_jx_0+a_j\}\in(0,1)$.  Define the finite local
numerator
\begin{align}
 B_i^{\boldsymbol r}(\boldsymbol\lambda)
 ={}&(-1)^{r_i}(1-\lambda_i)_{r_i}
          p_i^{\lambda_i-r_i-1}\notag\\
 &\quad\cdot\sum_{x_0\in S_i}[t^{r_i}]
   \prod_{j\ne i}\partial_y^{r_j}
          \mathscr G(\lambda_j,\eta_{ij}(x_0)+p_jt).
 \label{xid:eq:triple-local-B}
\end{align}
It is holomorphic near the spectral origin.  It uses a finite Taylor
polynomial and values of Stieltjes functions and their argument
derivatives at nonsingular points.

\begin{theorem}[All-index triple finite-part identity]
\label{xid:thm:triple-FP}
The function
\begin{equation}\label{xid:eq:triple-holomorphic-J}
 \mathcal J_{\boldsymbol r}(\boldsymbol\lambda)
 =\partial_{\boldsymbol a}^{\boldsymbol r}
       \mathcal L(\boldsymbol\lambda)
       -\sum_{i=1}^3\frac{B_i^{\boldsymbol r}(\boldsymbol\lambda)}{\lambda_i}
\end{equation}
is holomorphic in a sufficiently small neighborhood of the origin.
For all $m_1,m_2,m_3\ge0$,
\begin{align}
 &\FP_x\int_0^1\prod_{i=1}^3
       \gamma_{m_i}^{(r_i)}(\{p_ix+a_i\})\,dx\notag\\
 &\qquad=\left(\prod_i m_i!\right)
       [\lambda_1^{m_1}\lambda_2^{m_2}\lambda_3^{m_3}]
          \mathcal J_{\boldsymbol r}(\boldsymbol\lambda).
 \label{xid:eq:triple-FP-value}
\end{align}
Together with \eqref{xid:eq:triple-cones}, this is an explicit finite
colored-Tornheim, Gamma, and local Stieltjes coefficient formula for
every choice of the six indices.
\end{theorem}

\begin{proof}
Begin with $\Re\lambda_i>r_i+1$. Argument differentiation and
integration then commute, including at moving grid points: the singular
pieces vanish to the required orders and the regular pieces match across
each periodic endpoint, so the boundary terms cancel. Hence
$\partial_{\boldsymbol a}^{\boldsymbol r}\mathcal L$ is the
meromorphic integral of the product of the three argument-derivative
generators.

Near $x_0\in S_i$ the singular term of factor $i$ is
\[
 (-1)^{r_i}(1-\lambda_i)_{r_i}
 p_i^{\lambda_i-r_i-1}t^{\lambda_i-r_i-1}.
\]
The other two factors are smooth there.  Subtract their Taylor
polynomial through degree $r_i$.  The degree-$r_i$ term alone has
denominator $\lambda_i$ near the origin; summed over $S_i$ its
numerator is exactly \eqref{xid:eq:triple-local-B}.
As in Theorem~\ref{xid:thm:coordinate-monomial}, its meromorphic integral
over $(0,h)$ is $B_i h^{\lambda_i}/\lambda_i$, whereas the generating
function of its coordinate finite parts is
$B_i(h^{\lambda_i}-1)/\lambda_i$.  Their difference is the complete
quantity $B_i/\lambda_i$.

The lower Taylor degrees have denominators bounded away from zero
on a small spectral polydisc.  The remainders are locally normally
integrable there.  A partition of unity treats the finitely many
disjoint singular points independently.  Subtracting all three
complete numerators therefore leaves a holomorphic function whose
Taylor coefficients are exactly the stated finite-part integrals.
The factorials in \eqref{xid:eq:triple-FP-value} follow from the exponential
Stieltjes generating convention.
\end{proof}

For zero argument derivatives the local term simplifies to
\[
 B_i^{\boldsymbol0}(\boldsymbol\lambda)
 =p_i^{\lambda_i-1}\sum_{x_0\in S_i}
             \prod_{j\ne i}\mathscr G(\lambda_j,\eta_{ij}(x_0)).
\]
Even here its full dependence on $\lambda_i$ must be retained; using
only $B_i(0,\lambda_j,\lambda_k)$ omits the powers of $\log p_i$.
Thus the character calculation and the ambient-coordinate theorem are
separate necessary parts of the reduction.

\subsection{Independent polynomial checks}

At positive integer $\lambda_i$, the identity
$\zeta(1-\lambda_i,y)=-B_{\lambda_i}(y)/\lambda_i$ makes
\eqref{xid:eq:triple-K} a rational piecewise-polynomial integral whenever
the shifts are rational.  This provides a check independent of Fourier
or special-function quadrature.  For
\[
 (p_1,p_2,p_3)=(1,2,3),\qquad
 (a_1,a_2,a_3)=(0,1/5,1/7),
\]
the exact integrations give
\begin{align}
 K(2,2,2)&=-\frac{49960598339}{1715322420000000},\notag\\
 K(2,3,2)&=-\frac{1641301145317}{315190494675000000}.
 \label{xid:eq:triple-rational-checks}
\end{align}
The supplied program compares these exact rational values with the
independent filtered Mellin integrals \eqref{xid:eq:weighted-T-mellin}.
The spectral check does not replace the local-subtraction proof for
Stieltjes jets, which is given above.

% END sections/triple_dilation.tex


% BEGIN sections/fully_shifted_height_one.tex
% Research section prepared by harmonic_jets agent.
% Requires amsmath, amssymb, amsthm; theorem/corollary/proposition/remark envs.
% Citation keys proposed in harmonic_sources.bib (or adapt to main bibliography).
\section{Fully shifted height-one sums: an entire difference and exact Laurent data}
\label{xid:hsh:sec:main}

The inner and outer shifts of a nested harmonic sum are different pieces
of data. The earlier shifted height-one families in this project move the
outer denominator while retaining ordinary inner harmonic numbers. The
newest nested-jet report instead gives all nested denominators a common
variable exponent. Here every denominator has the same shift, while only
the outer exponent varies. A Gauss connection formula shows that changing
the common shift preserves a universal singular part, up to an explicit
Gamma factor. The remaining function is entire in the spectral variable,
and its values at every nonpositive integer are finite Stirling transforms.

At depth one this family belongs to the harmonic Hurwitz-zeta literature
of Karg\i n, Dil, Cenkci and Can \cite{hsh:Kargin2023}; the recent preprint
of Kara \"Ozt\"urk and Can treats the two-parameter depth-one family and
negative Laurent coefficients \cite{hsh:KaraCan2026}. The latter comparison
is based on its retrieved abstract: its complete text was not available
in this research environment. Thus the depth-one meromorphic continuation
and negative-integer evaluation problem are not claimed as new. The
all-depth entire difference, finite Stirling rule, and complete pole-order
classification below were derived here and were not found in the
inspected repository material; priority over all literature is not asserted.

\subsection{The common-shift convention}

Let $\Re a>0$. All powers of $n+a$ use the principal logarithm, and $u$
will remain in a sufficiently small disc centered at zero. Set
\begin{align}
 e_r(n;a)&=[u^r]\prod_{j=0}^{n-1}\left(1+\frac{u}{j+a}\right),
       &e_0(n;a)&=1,\label{xid:hsh:eq:er}\\
 M_r(s;a)&=\sum_{n=0}^{\infty}\frac{e_r(n;a)}{(n+a)^s},
       &&\Re s>1,\label{xid:hsh:eq:Mr}\\
 F_a(s,u)&=\sum_{n=0}^{\infty}
       \frac{(a+u)_n}{(a)_n(n+a)^s}
       =\sum_{r\ge0}M_r(s;a)u^r.
       \label{xid:hsh:eq:F}
\end{align}
The first product is empty when $n=0$. Equivalently,
\[
 M_r(s;a)=\sum_{n_0>n_1>\cdots>n_r\ge0}
       \frac{1}{(n_0+a)^s(n_1+a)\cdots(n_r+a)}.
\]
Thus $M_0(s;a)=\zeta(s,a)$ and $M_r(s;1)=\zeta(s,\{1\}^r)$.
If $H_n^{(k)}(a)=\sum_{j=0}^{n-1}(j+a)^{-k}$, then
\[
 e_1(n;a)=H_n^{(1)}(a),\qquad
 e_2(n;a)=\frac{(H_n^{(1)}(a))^2-H_n^{(2)}(a)}2.
\]
The general $e_r$ is the corresponding elementary-symmetric Bell
polynomial. Uniformly on compact $a$-sets,
$e_r(n;a)=O((1+\log n)^r)$; hence every fixed-depth series converges
normally for $\Re s>1$. The Gamma-ratio Dirichlet series in \eqref{xid:hsh:eq:F} converges normally
when $\Re s>1+\Re u$, locally uniformly in its parameters. Interchange
with its depth Taylor series is justified, for example, on
$\Re s>1+|u|$; at each fixed $\Re s>1$ one may choose a sufficiently
small disc in $u$. The direct-series estimate follows from
\[
 \frac{(a+u)_n}{(a)_n}
 =\frac{\Gamma(a)}{\Gamma(a+u)}
    \frac{\Gamma(n+a+u)}{\Gamma(n+a)}
 =O(n^{\Re u}).
\]
Gamma-ratio asymptotics and their uniform sector versions are classical
\cite{hsh:DLMFGamma}.

Introduce
\begin{align}
 C_a(u)&=\frac{\Gamma(a)}{\Gamma(a+u)}
        =\sum_{j\ge0}c_j(a)u^j,\label{xid:hsh:eq:Ca}\\
 K_a(u)&=\Gamma(1+u)C_a(u)
        =\sum_{j\ge0}k_j(a)u^j.\label{xid:hsh:eq:Ka}
\end{align}
The Gamma logarithm is analytic on the right half-plane, so
\begin{equation}\label{xid:hsh:eq:Cexp}
 C_a(u)=\exp\left(-\sum_{j\ge1}
                \frac{\psi^{(j-1)}(a)}{j!}u^j\right).
\end{equation}
Consequently every coefficient is a finite polynomial in polygamma
functions. In particular, with $P=\psi(a)$, $Q=\psi'(a)$, and
$R=\psi''(a)$,
\[
 c_0=1,\quad c_1=-P,\quad c_2=\frac{P^2-Q}{2},\quad
 c_3=\frac{-P^3+3PQ-R}{6}.
\]
The standard polygamma--Hurwitz identities can convert these coordinates
to zeta values when useful \cite{hsh:DLMFPolygamma,DLMFhurwitz}.

\subsection{The exact entire difference}

\begin{theorem}[Universal singular part under a common shift]
\label{xid:hsh:thm:entire}
For $\Re a>0$ and $u$ in a sufficiently small disc centered at zero,
\begin{equation}\label{xid:hsh:eq:D}
 D(s,u;a):=F_a(s,u)-K_a(u)F_1(s,u)
\end{equation}
extends jointly holomorphically to all $s\in\mathbb C$, locally in $a,u$.
Consequently, for every $r\ge0$,
\begin{equation}\label{xid:hsh:eq:Dr}
 D_r(s;a):=M_r(s;a)-\sum_{j=0}^rk_j(a)M_{r-j}(s;1)
\end{equation}
is entire in $s$.

For every integer $m\ge0$, its complete depth-generating value is
\begin{equation}\label{xid:hsh:eq:T}
 \boxed{\quad D(-m,u;a)
 =-\sum_{j=1}^{m+1}
       {\left\{\begin{matrix}m+1\\j\end{matrix}\right\}}
       \frac{(a-1)^{\underline j}}{u+j}.\quad}
\end{equation}
Here braces denote Stirling numbers of the second kind and
$x^{\underline j}=x(x-1)\cdots(x-j+1)$.
Equivalently,
\begin{equation}\label{xid:hsh:eq:Delta}
 \boxed{\quad
 D_r(-m;a)=\Delta_{m,r}(a):=
 (-1)^{r+1}\sum_{j=1}^{m+1}
       {\left\{\begin{matrix}m+1\\j\end{matrix}\right\}}
       \frac{(a-1)^{\underline j}}{j^{r+1}}.
 \quad}
\end{equation}
Thus every value in \eqref{xid:hsh:eq:Delta} is a rational polynomial in $a$
of degree at most $m+1$.
\end{theorem}

\begin{proof}
In the common half-plane of convergence, the Mellin kernel of $F_a$ is
\begin{equation}\label{xid:hsh:eq:Mellin}
 \Gamma(s)F_a(s,u)=\int_0^{\infty}t^{s-1}G_a(t,u)\,dt,
 \qquad G_a(t,u)=e^{-at}
          {}_2F_1(1,a+u;a;e^{-t}).
\end{equation}
Termwise integration follows from absolute convergence. The Gauss
connection formula at $z=1$ \cite{hsh:DLMFGauss} simplifies here to
\begin{equation}\label{xid:hsh:eq:kernelconnection}
 \begin{split}
 G_a(t,u)&=K_a(u)\frac{e^{-t}}{(1-e^{-t})^{1+u}}+R_a(t,u),\\
 R_a(t,u)&=-\frac{a-1}{1+u}e^{-at}
       {}_2F_1(1,a+u;2+u;1-e^{-t}).
 \end{split}
\end{equation}
Indeed its second local hypergeometric factor is
${}_2F_1(a-1,-u;-u;1-z)=z^{1-a}$, and its two connection
coefficients reduce to $K_a(u)$ and $-(a-1)/(1+u)$.
The identity is first read for noninteger $u$, where the generic
connection formula applies, and then continued through $u=0$.
The simplified two terms are individually analytic in $u$ near zero.

The function $R_a(t,u)$ is analytic in $t$ near zero, jointly in the
other parameters. Both $G_a$ and the first term in
\eqref{xid:hsh:eq:kernelconnection} decay exponentially as $t\to+\infty$,
locally uniformly for $\Re a>0$. Thus the same holds for $R_a$.
Write $R_a(t,u)=\sum_{j\ge0}r_j(a,u)t^j$. For any $L\ge1$,
\[
 \begin{split}
 \Gamma(s)D(s,u;a)={}&\sum_{j=0}^{L-1}\frac{r_j(a,u)}{s+j}\\
 &+\int_0^1t^{s-1}\left(R_a(t,u)-\sum_{j=0}^{L-1}r_j(a,u)t^j\right)dt
  +\int_1^{\infty}t^{s-1}R_a(t,u)\,dt.
 \end{split}
\]
The remainder integrals are jointly holomorphic for $\Re s>-L$.
All their parameter derivatives have the same local integrability,
possibly with fixed powers of $\log t$. Each displayed simple pole
at $s=-j$ is canceled by the simple zero of $1/\Gamma(s)$ there.
Increasing $L$ proves joint entireness in $s$ and gives
\begin{equation}\label{xid:hsh:eq:Dlocal}
 D(-m,u;a)=(-1)^m m!\,[t^m]R_a(t,u).
\end{equation}
For fixed $m$, the right side is a polynomial in $a$ of degree at most
$m+1$: only the first $m+1$ terms of the local hypergeometric series
can contribute, and the $j$-th term has a factor $(1-e^{-t})^j$.
Its coefficients are rational functions of $u$, analytic near zero.

It remains to identify this polynomial. For a positive integer $A$,
shifting the index of the defining series gives the exact finite
identity
\[
 D(s,u;A)=-K_A(u)\sum_{k=1}^{A-1}
           \frac{(1+u)_{k-1}}{(k-1)!\,k^s}.
\]
At $s=-m$, use
\[
 k^m=\sum_{j=1}^{m+1}
       {\left\{\begin{matrix}m+1\\j\end{matrix}\right\}}
         (k-1)^{\underline{j-1}}.
\]
The finite binomial summation yields
\[
 K_A(u)\sum_{k=1}^{A-1}
       \frac{(1+u)_{k-1}}{(k-1)!}
              (k-1)^{\underline{j-1}}
       =\frac{(A-1)^{\underline j}}{u+j}.
\]
For completeness, after putting $\ell=k-j$, the nonzero summands are
$(1+u)_{j-1}(u+j)_\ell/\ell!$; their sum is
$(1+u)_{j-1}(u+j+1)_{A-j-1}/(A-j-1)!$.
Multiplication by $K_A=(A-1)!/(1+u)_{A-1}$ gives the stated quotient.
When $A\le j$, both sides are zero. Therefore \eqref{xid:hsh:eq:T} holds
at infinitely many positive integers $A$. Both sides are polynomials
in $a$ of bounded degree, so they agree identically. Finally expand
$(u+j)^{-1}=\sum_{r\ge0}(-1)^ru^r/j^{r+1}$ to obtain
\eqref{xid:hsh:eq:Delta}.
\end{proof}

At $r=0$, the theorem reduces to the usual Bernoulli-polynomial
identity for $\zeta(-m,a)-\zeta(-m)$. At positive depths it gives a
finite rational interpolation of the entire difference after the
Gamma mixing of depths. The theorem does not identify
$M_r(-m;a)$ with a convergent series: the defining series diverges
there, and the statement concerns its meromorphic continuation.

\subsection{A bridge to independent inner and outer shifts}

The same argument unifies the earlier outer-shift calculus with the
common-shift family. Let $\Re a>0$, $\Re b>0$, and
$\Re c>0$, where $c=1+b-a$. Define
\[
 F_{a,b}(s,u)=\sum_{n\ge0}\frac{(a+u)_n}{(a)_n(n+b)^s},
 \qquad B_c(s,u)=\sum_{n\ge0}\frac{(1+u)_n}{n!(n+c)^s}.
\]
The function $B_c$ is precisely the previous generator with an outer
shift only. At $a=1$, $F_{a,b}=B_b$; at $b=a$, $F_{a,b}=F_a$.

\begin{corollary}[Two shifts and a finite binomial correction]
\label{xid:hsh:cor:twoshifts}
The difference
\[
 D_{a,b}(s,u)=F_{a,b}(s,u)-K_a(u)B_{1+b-a}(s,u)
\]
is entire in $s$, locally in its other parameters. For $m\ge0$,
\begin{equation}\label{xid:hsh:eq:twoshifts}
 D_{a,b}(-m,u)
   =\sum_{v=0}^m\binom mv(b-a)^{m-v}D(-v,u;a).
\end{equation}
Every term on the right is the explicit finite Stirling sum in
\eqref{xid:hsh:eq:T}. In depth $r$ the corresponding correction is
\[
 [u^r]D_{a,b}(-m,u)
    =\sum_{v=0}^m\binom mv(b-a)^{m-v}\Delta_{v,r}(a).
\]
\end{corollary}
\begin{proof}
Multiplication of the Mellin kernel in \eqref{xid:hsh:eq:kernelconnection}
by $e^{-(b-a)t}$ gives
\[
 e^{-bt}{}_2F_1(1,a+u;a;e^{-t})
   =K_a(u)\frac{e^{-(1+b-a)t}}{(1-e^{-t})^{1+u}}
         +e^{-(b-a)t}R_a(t,u).
\]
The last kernel is analytic at zero and decays exponentially at infinity
under the stated hypotheses. The same Mellin subtraction proves
entireness. Taking its $m$-th Taylor coefficient and applying
\eqref{xid:hsh:eq:Dlocal} proves \eqref{xid:hsh:eq:twoshifts} by the ordinary
binomial theorem.
\end{proof}

Thus no new singular function is required to separate arbitrary inner
and outer shifts in this domain: the existing outer-shift generator
carries the singular terms and a fully explicit polynomial correction
accounts for the values of the entire remainder at nonpositive integers.
The higher spectral derivatives of this remainder remain a separate
identity problem.

\subsection{The complete Laurent principal part and finite part}

Define the rational polynomials $b_j(u)$ and $W_m(u)$ by
\begin{align}
 e^{-t}\left(\frac{t}{1-e^{-t}}\right)^{1+u}
       &=\sum_{j\ge0}b_j(u)t^j,\label{xid:hsh:eq:b}\\
 W_m(u)&=(u-m)_{m+1}b_{m+1}(u).
       \label{xid:hsh:eq:W}
\end{align}
The logarithm of the parenthesized factor is the analytic germ at zero
whose value there is zero. Each $W_m$ is a rational polynomial divisible
by $u$. The first three are
\begin{align}
 W_0(u)&=\frac{u(u-1)}2,\label{xid:hsh:eq:W0}\\
 W_1(u)&=\frac{u(u-1)(u-2)(3u-1)}{24},\\
 W_2(u)&=\frac{u^2(u-1)^2(u-2)(u-3)}{48}.
\end{align}

\begin{theorem}[All depths, every nonpositive integer]
\label{xid:hsh:thm:Laurent}
Let $m,r\ge0$ and $\Re a>0$. Then
\begin{equation}\label{xid:hsh:eq:Laurent}
 \boxed{\quad
 M_r(-m+\varepsilon;a)=
 \sum_{h=0}^r
     \frac{[u^{r+1-h}]\{C_a(u)W_m(u)\}}{\varepsilon^h}
   +\Delta_{m,r}(a)+O(\varepsilon).
 \quad}
\end{equation}
In particular,
\begin{equation}\label{xid:hsh:eq:FP}
 \operatorname{FP}_{s=-m}M_r(s;a)
  =[u^{r+1}]\{C_a(u)W_m(u)\}+\Delta_{m,r}(a).
\end{equation}
The finite part belongs to
$\mathbb Q[a,\psi(a),\psi'(a),\ldots,\psi^{(r-1)}(a)]$,
where at $r=0$ the polygamma list is empty.
\end{theorem}

\begin{proof}
For the universal unshifted kernel, Mellin subtraction gives, near
$(s,u)=(-m,0)$,
\[
 F_1(-m+\varepsilon,u)
 =\frac{b_{m+1}(u)}{\Gamma(-m+\varepsilon)(\varepsilon-u)}
       +\frac{Q_m(\varepsilon,u)}{\Gamma(-m+\varepsilon)},
\]
where $Q_m$ is jointly holomorphic. The last term contributes only
$O(\varepsilon)$ after taking any fixed coefficient in $u$.
Expand $1/\Gamma(-m+\varepsilon)$ in powers of $\varepsilon$;
its constant coefficient is zero. From the first term, the coefficient
of $u^r\varepsilon^{-h}$, for $h\ge0$, is
\[
 [u^{r+1-h}]\frac{K_a(u)b_{m+1}(u)}{\Gamma(u-m)}
   =[u^{r+1-h}]\{C_a(u)W_m(u)\}.
\]
The second equality is the Gamma recurrence. Finally add the entire
function $D(s,u;a)$ and apply \eqref{xid:hsh:eq:Delta}. Since $W_m$ has
no constant coefficient, at most the first $r$ derivatives of
$\log\Gamma(a+u)$ can occur, giving the stated algebra.
\end{proof}

The polynomial correction in \eqref{xid:hsh:eq:FP} is essential. For instance,
with $r=m=0$, omitting it gives $-1/2$, whereas the correct value is
$\zeta(0,a)=1/2-a$. Thus the singular-coefficient rule for the earlier
outer-shift family cannot be transferred to the fully shifted family
without its analytic-kernel contribution. This is an extension guard,
not a claim that the earlier report made this incorrect transfer.

\begin{corollary}[Exact pole orders]
\label{xid:hsh:cor:poles}
At $s=1$, the order of the pole of $M_r$ is $r+1$.
At $s=0$ and every negative odd integer $s=-m$, the order is exactly $r$
for $r\ge1$, with leading coefficient $\zeta(-m)$.
At a negative even integer $s=-m$, $m\ge2$, the order is exactly $r-1$
for $r\ge2$, with leading coefficient
\begin{equation}\label{xid:hsh:eq:evenleading}
 \frac{m+1}{2m}B_m.
\end{equation}
At these negative even integers, $M_1$ is holomorphic.
All the specified leading coefficients are independent of $a$.
\end{corollary}

\begin{proof}
The assertion at one follows from the next theorem. From
\eqref{xid:hsh:eq:b},
\[
 [u]W_m(u)=(-1)^m m!\,b_{m+1}(0)=\zeta(-m).
\]
This is nonzero for $m=0$ and for positive odd $m$. If $m\ge2$ is even,
$b_{m+1}(0)=B_{m+1}/(m+1)!=0$. The two elementary expansions
\[
 \frac{t}{e^t-1}=1-\frac t2+
       \sum_{q\ge1}\frac{B_{2q}}{(2q)!}t^{2q},\qquad
 \log\frac{t}{1-e^{-t}}
       =\frac t2-\sum_{q\ge1}
          \frac{B_{2q}}{2q(2q)!}t^{2q}
\]
give
\[
 [t^{m+1}]\frac{t}{e^t-1}\log\frac{t}{1-e^{-t}}
       =\frac{m+1}{2m}\frac{B_m}{m!}.
\]
Multiplication by the leading term $m!u$ of $(u-m)_{m+1}$ proves
$[u^2]W_m=(m+1)B_m/(2m)$. This is nonzero, since every positive even
Bernoulli number is nonzero \cite{hsh:DLMFBernoulli}. Thus the zero of
$W_m$ at zero is exactly double in this case and simple in the preceding
case. Since $C_a(0)=1$ and $D$ is entire, the pole orders and leading
coefficients follow directly from \eqref{xid:hsh:eq:Laurent}.
\end{proof}

\subsection{The principal Stieltjes point}

\begin{theorem}[Gamma coefficients at the principal pole]
\label{xid:hsh:thm:one}
For $r\ge0$,
\begin{equation}\label{xid:hsh:eq:one}
 M_r(1+\varepsilon;a)
 =\sum_{j=0}^r\frac{c_j(a)}{\varepsilon^{r+1-j}}
       +c_{r+1}(a)+O(\varepsilon).
\end{equation}
Also
\begin{equation}\label{xid:hsh:eq:D1}
 D(1,u;a)=\frac{K_a(u)-1}{u},\qquad D_r(1;a)=k_{r+1}(a).
\end{equation}
The apparent singularity on the right is removable.
\end{theorem}

\begin{proof}
Let $p_n=(a+u)_n/(a)_n$. Its exact difference is
$p_{n+1}-p_n=up_n/(n+a)$. For small $\Re u<0$, $p_n\to0$;
therefore telescoping proves $F_a(1,u)=-1/u$.
The leading kernel term in \eqref{xid:hsh:eq:kernelconnection} shows that
\[
 F_a(1+\varepsilon,u)=\frac{C_a(u)}{\varepsilon-u}
                         +Q_a(\varepsilon,u)
\]
with $Q_a$ jointly holomorphic near zero. Evaluate at $\varepsilon=0$
for $\Re u<0$ to obtain $Q_a(0,u)=(C_a(u)-1)/u$, and continue
holomorphically through $u=0$. Coefficient extraction proves
\eqref{xid:hsh:eq:one}. Applying the same telescoping value to the
entire difference gives \eqref{xid:hsh:eq:D1}.
\end{proof}

At depth one, the commonly used inclusive harmonic Hurwitz function is
\[
 \zeta_H(s,a)=\sum_{n\ge0}\frac{\sum_{j=0}^{n}(j+a)^{-1}}{(n+a)^s}
            =M_1(s;a)+\zeta(s+1,a).
\]
Thus \eqref{xid:hsh:eq:one} gives
\[
 \zeta_H(1+\varepsilon,a)
 =\frac1{\varepsilon^2}-\frac{\psi(a)}{\varepsilon}
       +\frac{\psi(a)^2+\psi'(a)}2+O(\varepsilon).
\]
The strict and inclusive conventions differ even at the finite part.
The formula at all depths gives the Gamma-polygamma coefficients of the
strict height-one family; higher regular Laurent coefficients are not
claimed to reduce to this same finite algebra.

\subsection{Explicit values and Gamma antiderivatives}

Writing $P=\psi(a)$ and $Q=\psi'(a)$, the first examples are
\begin{align}
 M_1(\varepsilon;a)
  &=-\frac1{2\varepsilon}+a-\frac12+\frac P2+O(\varepsilon),
       \label{xid:hsh:eq:M10}\\
 M_1(-1+\varepsilon;a)
  &=-\frac1{12\varepsilon}
       +\frac{a^2+a}{4}-\frac18+\frac P{12}+O(\varepsilon),\\
 M_1(-2;a)&=\frac{8a^3+6a^2-2a-3}{72},
       \label{xid:hsh:eq:M12}\\
 M_2(\varepsilon;a)
  &=-\frac1{2\varepsilon^2}+\frac{1+P}{2\varepsilon}
       +1-a-\frac P2-\frac{P^2-Q}{4}+O(\varepsilon),\\
 M_2(-1+\varepsilon;a)
  &=-\frac1{12\varepsilon^2}
       +\frac{3/8+P/12}{\varepsilon}
       +\frac{-3a^2-15a+8-9P-P^2+Q}{24}
       +O(\varepsilon).\label{xid:hsh:eq:M21}
\end{align}
At the half shift, put $L=\gamma+2\log2$, so
$\psi(1/2)=-L$ and $\psi'(1/2)=\pi^2/2$. Then
\begin{align}
 M_1(-2;1/2)&=-\frac1{48},\\
 M_2(-2+\varepsilon;1/2)
       &=\frac1{8\varepsilon}+\frac L8-\frac{19}{288}
                    +O(\varepsilon).
\end{align}
In the convergent half-plane, for example,
$M_2(s;1/2)$ is one half the Dirichlet series whose numerator is
$(\sum_{j<n}(j+1/2)^{-1})^2-\sum_{j<n}(j+1/2)^{-2}$.
Its displayed pole and finite part follow without a numerical relation search.

There is an exact antiderivative for every depth-one finite part. Put
\[
 P_m(a)=[u^2]W_m(u)+
       \sum_{j=1}^{m+1}{\left\{\begin{matrix}m+1\\j\end{matrix}\right\}}
                 \frac{(a-1)^{\underline j}}{j^2}\in\mathbb Q[a].
\]
Then
\begin{equation}\label{xid:hsh:eq:depthonegeneral}
 \operatorname{FP}_{s=-m}M_1(s;a)
       =P_m(a)-\zeta(-m)\psi(a),
\end{equation}
and on the right half-plane
\begin{equation}\label{xid:hsh:eq:primitive}
 \int\operatorname{FP}_{s=-m}M_1(s;a)\,da
       =\mathcal P_m(a)-\zeta(-m)\log\Gamma(a)+\text{constant},
 \qquad \mathcal P_m'=P_m.
\end{equation}
For even $m\ge2$, the Gamma term vanishes and the value in
\eqref{xid:hsh:eq:depthonegeneral} is itself a polynomial. For $m=0,1$,
\begin{align*}
 \int\operatorname{FP}_{s=0}M_1(s;a)\,da
       &=\frac{a^2-a}{2}+\frac12\log\Gamma(a)+\text{constant},\\
 \int\operatorname{FP}_{s=-1}M_1(s;a)\,da
       &=\frac{a^3}{12}+\frac{a^2}{8}-\frac a8
                      +\frac1{12}\log\Gamma(a)+\text{constant}.
\end{align*}
Higher derivatives in $a$ consequently give explicit polygamma formulas.
This assertion is about finite Laurent coefficients; it does not replace
a divergent defining series by termwise integration at a negative order.

\subsection{Reproduction and further questions}

The accompanying verification script performs $72$ floating-point checks
at $75$ decimal working digits. Its maximum absolute residual is below
$8.5\times10^{-52}$. Thirty checks compare \eqref{xid:hsh:eq:T} with an
independent continuation of the original Gamma-ratio Dirichlet series,
using its Bernoulli-polynomial asymptotic coefficients and Hurwitz-zeta
tails. Two repeats with a larger cutoff and more tail terms have residuals
below $1.7\times10^{-64}$. The remaining checks audit the Gauss connection,
its mixed Taylor coefficients, and \eqref{xid:hsh:eq:D1}. These are
floating-point checks, not interval certificates; the all-parameter
proof is the analytic argument above.

A constructive next step is to determine the higher regular Taylor
coefficients of the entire difference $D_r(s;a)$ at $s=-m$. Unlike its
value, its first derivative includes a genuine integral of the analytic
kernel and need not be a rational polynomial. A second direction is to
retain independent perturbations of the inner exponents: the finite-product
model then requires more than one spectral direction, and the singular
part need not be governed by a single Gamma quotient. A third direction is to compare the finite
Stirling transforms in \eqref{xid:hsh:eq:Delta} with established
poly-Bernoulli and generalized Bernoulli polynomial systems, with exact
normalization matching before claiming an identification. Finally,
character-weighted sums in the common shift may yield further finite
cyclotomic reductions of the polygamma algebra in \eqref{xid:hsh:eq:FP}.

% END sections/fully_shifted_height_one.tex


% BEGIN sections/weighted_arctangent.tex
\section{A weighted arctangent integral in classical trilogarithms}
\label{xid:wat:sec}

The canonical manuscript, in the paragraph following
\texttt{cleo:eq:lintanh}, considers the weighted analogue of the
Reshetnikov--Lee bilinear integral,
\begin{equation}\label{xid:wat:eq:K}
 K(p,q)=\int_0^{\pi/2}
 \frac{\arctan(p\sin x)\arctan(q\sin x)}{\sin x}\,dx,
 \qquad p,q\in\mathbb R.
\end{equation}
It explains why the unweighted reduction to one Legendre chi function
does not carry over, but gives no finite evaluation for this family.
We supply a two-term formula in a single even function built from
classical dilogarithms and trilogarithms at real arguments in $(0,1]$.
We also give a normalized ordinary antiderivative and an
arbitrary-order Euler integration ladder. The proof of the first
result uses parameter differentiation and hyperbolic addition;
the original unweighted proof of Lee~\cite{wat:Lee} also uses
hyperbolic parameters. The elementary logarithmic reduction used
later is part of standard weight-three polylogarithm calculus.
Its underlying two-logarithm kernel was treated for real parameters
by David H.~\cite{wat:DavidH}; the polarization method is not new.
The contribution relative to the inspected repository is the explicit
evaluation of its weighted two-parameter family, with real domains,
normalizations, and branch prescriptions proved. Historical priority
for this exact weighted evaluation has not been established.

The integral in \eqref{xid:wat:eq:K} is ordinary and convergent. Its integrand
is $pq\,x+O(x^3)$ at zero. The principal real arctangent is used
throughout. Principal complex logarithms and classical polylogarithms
are needed only for the more general ordinary primitive. Their series and integral
conventions agree with DLMF~\S25.12~\cite{wat:DLMF}.

\subsection{The hyperbolic addition formula}

Write
\[
 \chi_j(z)=\frac{\operatorname{Li}_j(z)-\operatorname{Li}_j(-z)}2
          =\sum_{k=0}^{\infty}\frac{z^{2k+1}}{(2k+1)^j}
 \qquad (|z|<1)
\]
for the Legendre chi function. The values used in the next theorem
are real and lie in its ordinary convergent series domain, including
the endpoints $\chi_2(1)=\pi^2/8$ and
$\chi_3(1)=7\zeta(3)/8$.

\begin{theorem}[A real two-term weighted evaluation]
\label{xid:wat:thm:main}
Define the even function
\begin{equation}\label{xid:wat:eq:Fclosed}
 \mathfrak F(s)=\frac{\pi^2|s|}{8}-\frac74\zeta(3)
    +|s|\chi_2(e^{-|s|})+2\chi_3(e^{-|s|}),
 \qquad s\in\mathbb R.
\end{equation}
Its value at zero is $\mathfrak F(0)=0$. The apparent lack of
smoothness from the absolute values is removable: $\mathfrak F$ is
real analytic on $\mathbb R$, and it is equivalently characterized by
\begin{equation}\label{xid:wat:eq:Fode}
 \mathfrak F''(s)=\frac{s}{2\sinh s},\qquad
 \mathfrak F''(0)=\frac12,\qquad
 \mathfrak F(0)=\mathfrak F'(0)=0.
\end{equation}
For every $p,q\in\mathbb R$, put
$\alpha=\operatorname{arcsinh}p$ and
$\beta=\operatorname{arcsinh}q$. Then
\begin{equation}\label{xid:wat:eq:main}
 \boxed{\displaystyle
 K(p,q)=\mathfrak F(\alpha+\beta)-\mathfrak F(\alpha-\beta).}
\end{equation}
In particular the formula includes both diagonals $p=\pm q$,
either zero parameter, and arbitrary signs without branch changes.
\end{theorem}

\begin{proof}
First, for $s>0$ use
$d\chi_j(e^{-s})/ds=-\chi_{j-1}(e^{-s})$ to obtain
\[
 \mathfrak F'(s)=\frac{\pi^2}{8}
              -\chi_2(e^{-s})-s\chi_1(e^{-s}),
 \qquad
 \mathfrak F''(s)=s\chi_0(e^{-s})
                =\frac{s}{2\sinh s}.
\]
The limits at zero follow from the endpoint values of $\chi_2$ and
$\chi_3$ and from
$s\chi_1(e^{-s})=\tfrac{s}{2}\log\coth(s/2)\to0$.
The unique twice-normalized primitive of the real analytic even
function $s/(2\sinh s)$ is real analytic and even. Hence
\eqref{xid:wat:eq:Fclosed} agrees with this primitive on the whole real
line; this proves all assertions in \eqref{xid:wat:eq:Fode}.

Differentiate \eqref{xid:wat:eq:K} once in each parameter. On compact
real parameter sets the differentiated integrand is bounded by an
integrable function, and therefore
\begin{align*}
 K_{pq}(p,q)
 &=\int_0^{\pi/2}
       \frac{\sin x\,dx}
       {(1+p^2\sin^2x)(1+q^2\sin^2x)}\\
 &=\frac{p\operatorname{arcsinh}p/\sqrt{1+p^2}
       -q\operatorname{arcsinh}q/\sqrt{1+q^2}}{p^2-q^2}
       \qquad(p^2\ne q^2).
\end{align*}
For the second equality use partial fractions and
\[
 \int_0^{\pi/2}\frac{\sin x}{1+p^2\sin^2x}\,dx
       =\frac{\operatorname{arcsinh}p}{p\sqrt{1+p^2}},
\]
whose zero-parameter value is $1$ and which follows by
$v=\cos x$.

Let $\widetilde K(\alpha,\beta)=K(\sinh\alpha,\sinh\beta)$.
Multiplication by $\cosh\alpha\cosh\beta$ and the addition formulas
for $\sinh$ now give
\begin{align}\label{xid:wat:eq:hyperbolicmixed}
 \widetilde K_{\alpha\beta}
 &=\frac{\alpha\sinh\alpha\cosh\beta
        -\beta\sinh\beta\cosh\alpha}
       {\sinh^2\alpha-\sinh^2\beta}\notag\\
 &=\frac12\left\{
      \frac{\alpha+\beta}{\sinh(\alpha+\beta)}
     +\frac{\alpha-\beta}{\sinh(\alpha-\beta)}\right\}.
\end{align}
Both quotients on the last line have their removable value $1$ at
zero. Thus the equality holds at $\alpha=\pm\beta$ by continuity.
By \eqref{xid:wat:eq:Fode}, this is the mixed derivative of
$\mathfrak F(\alpha+\beta)-\mathfrak F(\alpha-\beta)$.
Both this expression and $\widetilde K$ vanish on $\alpha=0$ and
$\beta=0$. Integration of their equal mixed derivatives over the
oriented rectangle from $(0,0)$ to $(\alpha,\beta)$ proves
\eqref{xid:wat:eq:main} for every real $\alpha,\beta$.
\end{proof}

For algebraic real $p,q$, the two arguments
$e^{-|\alpha+\beta|}$ and $e^{-|\alpha-\beta|}$ are algebraic.
Indeed $e^\alpha=p+\sqrt{1+p^2}$ and
$e^\beta=q+\sqrt{1+q^2}$ are positive algebraic numbers, and taking
the absolute value in either exponent selects an algebraic number
or its reciprocal. Thus \eqref{xid:wat:eq:main} expresses every such
integral by classical polylogarithms at real algebraic arguments and
logarithms of positive algebraic numbers. It asserts no independence
or minimality of the resulting constants.

The diagonal specializes particularly simply:
\begin{equation}\label{xid:wat:eq:Kdiagonal}
 K(p,p)=\mathfrak F(2\operatorname{arcsinh}p).
\end{equation}
For $\ell=\log(1+\sqrt2)$ and $r=3-2\sqrt2$, this gives the exact
identity
\begin{equation}\label{xid:wat:eq:K11}
 \boxed{\displaystyle
 \int_0^{\pi/2}\frac{\arctan^2(\sin x)}{\sin x}\,dx
 =\frac{\pi^2\ell}{4}-\frac74\zeta(3)
       +2\ell\chi_2(r)+2\chi_3(r).}
\end{equation}
Its numerical value is
$0.718064319741103407036328710693\ldots$; the decimal only
identifies the normalization, while the theorem proves the identity.

An exact weighted counterpart of the manuscript's Cleo integral
also follows. Set $\alpha=\log(1+\sqrt2)$ and
$\beta=\log\bigl((1+\sqrt5)/2\bigr)$. Then
\begin{align}\label{xid:wat:eq:weightedCleo}
 &\int_0^{\pi/2}\frac1{\sin x}
    \arctan^2\!\left(\frac{6\sin x}{3+\cos2x}\right)dx\notag\\
 &\qquad=\mathfrak F(2\alpha)
          +2\mathfrak F(\alpha+\beta)
          -2\mathfrak F(\alpha-\beta)
          +\mathfrak F(2\beta).
\end{align}
To verify the reduction put $s=\sin x\in[0,1]$. Since
$\arctan s+\arctan(s/2)<\pi/2$, the tangent addition formula gives
\[
 \arctan\frac{6\sin x}{3+\cos2x}
       =\arctan(\sin x)+\arctan(\tfrac12\sin x).
\]
Squaring and integrating yields
$K(1,1)+2K(1,1/2)+K(1/2,1/2)$, proving
\eqref{xid:wat:eq:weightedCleo} by \eqref{xid:wat:eq:main}.

\subsection{The weighted inverse-hyperbolic companion}

The same parameter method has a trigonometric form. It gives a
second exact family, including convergent improper endpoint values.

\begin{theorem}[An inverse-hyperbolic bilinear evaluation]
\label{xid:wat:thm:atanh}
For $|r|,|s|\leq1$ define
\[
 L(r,s)=\int_0^{\pi/2}
    \frac{\operatorname{arctanh}(r\sin x)
          \operatorname{arctanh}(s\sin x)}{\sin x}\,dx,
\]
using real inverse hyperbolic tangents and interpreting an endpoint
with $|r|=1$ or $|s|=1$ as an improper integral. Put
\begin{equation}\label{xid:wat:eq:Gclosed}
 \mathfrak G(t)=\frac74\zeta(3)
     -t\operatorname{Im}\chi_2(e^{it})
     -2\operatorname{Re}\chi_3(e^{it}),\qquad |t|\leq\pi.
\end{equation}
The chi values on the unit circle are defined by their absolutely
convergent series. Then $\mathfrak G$ is even, continuous on
$[-\pi,\pi]$, and real analytic on $(-\pi,\pi)$, with
\begin{equation}\label{xid:wat:eq:Gode}
 \mathfrak G''(t)=\frac{t}{2\sin t},\qquad
 \mathfrak G''(0)=\frac12,\qquad
 \mathfrak G(0)=\mathfrak G'(0)=0.
\end{equation}
For $a=\arcsin r$, $b=\arcsin s$, including either sign of either
parameter,
\begin{equation}\label{xid:wat:eq:atanhmain}
 \boxed{\displaystyle
 L(r,s)=\mathfrak G(a+b)-\mathfrak G(a-b).}
\end{equation}
In particular,
\begin{equation}\label{xid:wat:eq:atanh11}
 \boxed{\displaystyle
 \int_0^{\pi/2}\frac{\operatorname{arctanh}^2(\sin x)}{\sin x}\,dx
       =\frac72\zeta(3).}
\end{equation}
\end{theorem}

\begin{proof}
Conjugation of the chi series proves evenness and endpoint
continuity of \eqref{xid:wat:eq:Gclosed}. For $0<t<\pi$,
polylogarithm differentiation gives
\[
 \mathfrak G'(t)=\operatorname{Im}\chi_2(e^{it})
                    -t\operatorname{Re}\chi_1(e^{it}),
 \qquad
 \mathfrak G''(t)=t\operatorname{Im}\chi_0(e^{it})
                 =\frac{t}{2\sin t}.
\]
Here $\operatorname{Re}\chi_1(e^{it})=\tfrac12\log\cot(t/2)$,
so the limit of the first derivative at zero is zero. Also
$\mathfrak G(0)=0$ by $\chi_3(1)=7\zeta(3)/8$.
The normalized even primitive of $t/(2\sin t)$ is real analytic
on $(-\pi,\pi)$, proving \eqref{xid:wat:eq:Gode}.

For $|r|,|s|<1$, differentiation and partial fractions give
\[
 L_{rs}(r,s)=
 \frac{r\arcsin r/\sqrt{1-r^2}
       -s\arcsin s/\sqrt{1-s^2}}{r^2-s^2}
 \qquad(r^2\ne s^2).
\]
The underlying elementary integral is
$\int_0^{\pi/2}\sin x/(1-r^2\sin^2x)\,dx
 =\arcsin r/(r\sqrt{1-r^2})$, again obtained by $v=\cos x$.
For $\widetilde L(a,b)=L(\sin a,\sin b)$ the sine addition formulas
therefore yield
\[
 \widetilde L_{ab}(a,b)
   =\frac12\left\{
       \frac{a+b}{\sin(a+b)}+\frac{a-b}{\sin(a-b)}\right\},
 \qquad |a|,|b|<\pi/2.
\]
The zero quotients take their removable value $1$. This is the
mixed derivative of the right side of \eqref{xid:wat:eq:atanhmain},
and both sides vanish on the two axes. Rectangle integration proves
the identity in the open parameter square.

For all $|r|,|s|\leq1$ the absolute value of the integrand is bounded
by $\operatorname{arctanh}^2(\sin x)/\sin x$. This majorant is
$O(x)$ at zero and $O(\log^2(\pi/2-x))$ at the other endpoint,
so it is integrable. Dominated convergence extends the equality
to the closed square. Finally
$\chi_3(-1)=-7\zeta(3)/8$ and
$\operatorname{Im}\chi_2(-1)=0$ give
$\mathfrak G(\pi)=7\zeta(3)/2$, proving
\eqref{xid:wat:eq:atanh11}.
The endpoint value also follows directly from
$u=\operatorname{arctanh}(\sin x)$:
\[
 \int_0^\infty\frac{u^2}{\sinh u}\,du
    =4\sum_{k=0}^\infty\frac1{(2k+1)^3}
    =\frac72\zeta(3).
\]
\end{proof}

For algebraic $r,s\in[-1,1]$, the chi arguments in this theorem are
algebraic points on the unit circle, since
$e^{ia}=\sqrt{1-r^2}+ir$ and
$e^{ib}=\sqrt{1-s^2}+is$. The theorem has been proved directly over
real parameters; it does not depend on an unstated continuation of
the arctangent formula.

\subsection{The classical logarithmic kernel}

For $A,B>0$, we will need only the two types of kernel
\begin{equation}\label{xid:wat:eq:Ldef}
 \mathcal L(a,b)=\int_0^1
        \frac{\log(1-at)\log(1-bt)}{t}\,dt,
 \qquad (a,b)=(iA,iB)\text{ or }(iA,-iB).
\end{equation}
These integrals have no singularity on their integration path. Define
\begin{equation}\label{xid:wat:eq:P}
 P_c(y)=\log^2y\log(1-y/c)
       +2\log y\operatorname{Li}_2(y/c)
       -2\operatorname{Li}_3(y/c).
\end{equation}
The limiting convention is $P_1(1)=-2\zeta(3)$. In each other use of
$P_c(1)$ below, $1/c\notin[1,\infty)$, and thus
$P_c(1)=-2\operatorname{Li}_3(1/c)$ directly.

\begin{lemma}[A finite principal-branch logarithmic reduction]
\label{xid:wat:lem:L}
If $a=iA$, $b=iB$, $A>B>0$, or if $a=iA$, $b=-iB$, $A,B>0$, set
$r=(1-a)/(1-b)$. Then
\begin{equation}\label{xid:wat:eq:Lclosed}
 \boxed{\displaystyle
 \mathcal L(a,b)=\zeta(3)+\operatorname{Li}_3(b/a)
 +\frac{P_1(1-a)+P_1(1-b)-P_1(r)+P_{a/b}(r)}2.}
\end{equation}
For the same-sign case with $A<B$, interchange $a$ and $b$ before
applying this expression. At the diagonal,
\begin{equation}\label{xid:wat:eq:Ldiagonal}
 \mathcal L(a,a)=P_1(1-a)+2\zeta(3).
\end{equation}
All quantities in these prescriptions use principal branches; no
unstated continuation across a positive-real polylogarithm cut is needed.
\end{lemma}

\begin{proof}
Differentiation using
$d\operatorname{Li}_j(z)/dz=\operatorname{Li}_{j-1}(z)/z$ and
$\operatorname{Li}_1(z)=-\log(1-z)$ gives
\begin{equation}\label{xid:wat:eq:Pderivative}
 \frac{d}{dy}P_c(y)=\frac{\log^2y}{y-c}.
\end{equation}
The four extraneous terms cancel in pairs. Consequently
\[
 \int_0^1\frac{\log^2(1-at)}t\,dt
 =P_1(1-a)-P_1(1).
\]
For the cross term use the polarization identity
\[
 2\log(1-at)\log(1-bt)
 =\log^2(1-at)+\log^2(1-bt)
  -\log^2\frac{1-at}{1-bt}.
\]
Here the logarithm of the quotient equals the difference of the two
principal logarithms on the stated paths. Indeed, in the same-sign
case its argument is $-\arctan(At)+\arctan(Bt)\in(-\pi/2,0)$;
in the opposite-sign case it is
$-\arctan(At)-\arctan(Bt)\in(-\pi,0)$.

With $y=(1-at)/(1-bt)$ and $c=a/b$, elementary differentiation yields
\[
 \frac{dt}{t}=\left(\frac1{y-1}-\frac1{y-c}\right)dy.
\]
The quotient-log square therefore has integral
\[
 P_1(r)-P_1(1)-P_c(r)+P_c(1).
\]
Substitution into polarization proves \eqref{xid:wat:eq:Lclosed} and also
the diagonal formula.

For completeness, the branch check for the additional polylogarithms
is explicit. In the same-sign case $A>B$,
\[
 y(t)=\frac{1+ABt^2+i(B-A)t}{1+B^2t^2},
 \qquad
 \operatorname{Re}\frac{y(t)}c
 =\frac{B/A+B^2t^2}{1+B^2t^2}<1.
\]
Thus $y(t)$ is in the lower half-plane for $t>0$, while $y(t)/c$
starts in $(0,1)$ and avoids $[1,\infty)$.
In the opposite-sign case $y(t)$ is in the lower half-plane and
$y(t)/c$, with $c<0$, is in the upper half-plane, starting at
$-B/A<0$. These observations justify every use of
\eqref{xid:wat:eq:Pderivative} with the stated principal values. The
endpoint $y=1$ is removable in $P_1$ because
$\log^2y\log(1-y)\to0$ and
$\log y\operatorname{Li}_2(y)\to0$.
\end{proof}

Define the real symmetric kernel
\begin{equation}\label{xid:wat:eq:J}
 \mathcal J(A,B)=\frac12\operatorname{Re}
       \{\mathcal L(iA,-iB)-\mathcal L(iA,iB)\},
 \qquad A,B>0.
\end{equation}
Lemma~\ref{xid:wat:lem:L} makes this a finite expression in classical
trilogarithms, dilogarithms, logarithms, and $\zeta(3)$.
Conjugation of the logarithms gives the simpler integral meaning
\begin{equation}\label{xid:wat:eq:Jintegral}
 \mathcal J(A,B)=\int_0^1
              \frac{\arctan(At)\arctan(Bt)}t\,dt.
\end{equation}
To see the factor $1/2$, write
$\log(1-iAt)=\frac12\log(1+A^2t^2)-i\arctan(At)$;
the real part of the difference in \eqref{xid:wat:eq:J} is exactly twice
the product of arctangents.

\subsection{The ordinary primitive at a general endpoint}

\begin{theorem}[Finite classical evaluation and a normalized primitive]
\label{xid:wat:thm:primitive}
For $p,q\in\mathbb R$ put
\[
 a=e^{\operatorname{arcsinh}p}=p+\sqrt{1+p^2}>0,
 \qquad
 b=e^{\operatorname{arcsinh}q}=q+\sqrt{1+q^2}>0.
\]
For $0<\theta<\pi$ let $T=\tan(\theta/2)$. Then
\begin{align}\label{xid:wat:eq:primitive}
 &\int_0^\theta
   \frac{\arctan(p\sin x)\arctan(q\sin x)}{\sin x}\,dx
       \notag\\
 &\quad=\mathcal J(aT,bT)-\mathcal J(aT,T/b)
           -\mathcal J(T/a,bT)+\mathcal J(T/a,T/b).
\end{align}
In particular,
\begin{equation}\label{xid:wat:eq:fourkernel}
 \boxed{K(p,q)=\mathcal J(a,b)-\mathcal J(a,b^{-1})
       -\mathcal J(a^{-1},b)+\mathcal J(a^{-1},b^{-1}).}
\end{equation}
Together with Lemma~\ref{xid:wat:lem:L}, these are finite formulas using
only classical polylogarithms of orders at most three. They include
negative and zero values of $p,q$ and the diagonals $p=\pm q$.
\end{theorem}

\begin{proof}
The half-angle substitution $t=\tan(x/2)$ gives
\[
 \sin x=\frac{2t}{1+t^2},\qquad \frac{dx}{\sin x}=\frac{dt}{t}.
\]
Since $a-a^{-1}=2p$,
\begin{equation}\label{xid:wat:eq:atandifference}
 \arctan\frac{2pt}{1+t^2}
       =\arctan(at)-\arctan(t/a),\qquad t>0.
\end{equation}
The tangents of the two sides coincide. Both sides lie strictly
between $-\pi/2$ and $\pi/2$ because the two arctangents on the right
belong to $(0,\pi/2)$. Hence the equality has no additive multiple
of $\pi$, including when $p<0$.
Apply \eqref{xid:wat:eq:atandifference} for $p$ and $q$, expand the
product, and rescale each integral over $[0,T]$ to $[0,1]$.
Equation~\eqref{xid:wat:eq:Jintegral} gives \eqref{xid:wat:eq:primitive}.
The choice $\theta=\pi/2$ gives $T=1$ and proves
\eqref{xid:wat:eq:fourkernel}.
\end{proof}

The right side of \eqref{xid:wat:eq:primitive}, plus an arbitrary constant,
is an ordinary antiderivative in $\theta$ on $(0,\pi)$. Its constant
has been fixed by the limit zero at $\theta=0$. For algebraic $p,q$
and algebraic $T>0$, every argument in the finite formula is algebraic.
This includes rational multiples $\theta/\pi$ with $0<\theta<\pi$.
The statement is an explicit upper bound on the kinds of functions
needed for the expression; it asserts no independence or minimality
of its special values.

\subsection{An elementary differential certificate}

The elementary computation in the proof of
Theorem~\ref{xid:wat:thm:main} supplies a differential certificate for
both finite formulas, including their diagonal limits.

\begin{proposition}[Mixed derivative]
\label{xid:wat:prop:derivative}
For real $p,q$ with $p^2\ne q^2$,
\begin{equation}\label{xid:wat:eq:mixedderivative}
 \partial_p\partial_qK(p,q)
 =\frac{p\operatorname{arcsinh}p/\sqrt{1+p^2}
       -q\operatorname{arcsinh}q/\sqrt{1+q^2}}{p^2-q^2}.
\end{equation}
At $p^2=q^2$ the continuous value, for $p\ne0$, is
\begin{equation}\label{xid:wat:eq:mixeddiagonal}
 \frac1{2(1+p^2)}
 +\frac{\operatorname{arcsinh}p}{2p(1+p^2)^{3/2}},
\end{equation}
and the value at $(0,0)$ is $1$.
\end{proposition}

\begin{proof}
Equation~\eqref{xid:wat:eq:mixedderivative} was proved above. Take its
divided-difference limit by differentiating the numerator as a
function of $p^2$ to obtain \eqref{xid:wat:eq:mixeddiagonal}.
Its continuous limit at zero is $1$.
\end{proof}

Since $K(p,0)=K(0,q)=0$, its mixed derivative determines $K$ uniquely
by integration over the rectangle with corners $(0,0)$ and $(p,q)$.
This supplies an elementary differential certificate for any further
compressed trilogarithm representation of the same family.

\subsection{Every Euler integration order}

The half-angle reduction is not confined to weight three. For
integers $m\geq1$ set
\begin{equation}\label{xid:wat:eq:Lm}
 \mathcal L_m(a,b)=\frac1{(m-1)!}\int_0^1
    \frac{(-\log t)^{m-1}\log(1-at)\log(1-bt)}t\,dt.
\end{equation}
For the imaginary arguments used below this is an ordinary convergent
integral. The convention for double polylogarithms is
\[
 \operatorname{Li}_{r,1}(x,y)
 =\sum_{n>k\geq1}\frac{x^ny^k}{n^r k}.
\]

\begin{theorem}[A depth-two Euler primitive ladder]
\label{xid:wat:thm:ladder}
For $m\geq1$,
\begin{equation}\label{xid:wat:eq:LmMPL}
 \mathcal L_m(a,b)
 =\operatorname{Li}_{m+1,1}(a,b/a)
  +\operatorname{Li}_{m+1,1}(b,a/b).
\end{equation}
This is initially a series identity for $|a|,|b|<1$ and nonzero $a,b$.
For purely imaginary nonzero $a,b$, continue the two arguments
simultaneously along $(a,b)\mapsto(\lambda a,\lambda b)$,
$0<\lambda\leq1$. This radial prescription gives the value in
\eqref{xid:wat:eq:Lm}; no singularity is crossed.

Define
\[
 \mathcal J_m(A,B)=\frac12\operatorname{Re}
     \{\mathcal L_m(iA,-iB)-\mathcal L_m(iA,iB)\},
 \qquad A,B>0,
\]
and abbreviate
\[
 D_p(t)=\arctan(at)-\arctan(t/a)
       =\arctan\frac{2pt}{1+t^2}.
\]
Set
\begin{align}\label{xid:wat:eq:Imclosed}
 \mathcal I_m(T;p,q)
 ={}&\mathcal J_m(aT,bT)-\mathcal J_m(aT,T/b)\notag\\
   &-\mathcal J_m(T/a,bT)+\mathcal J_m(T/a,T/b).
\end{align}
Then
\begin{equation}\label{xid:wat:eq:Imint}
 \mathcal I_m(T;p,q)
 =\frac1{(m-1)!}\int_0^T
    \log^{m-1}(T/t)D_p(t)D_q(t)\,\frac{dt}{t}.
\end{equation}
Writing $\Theta_T=T\,d/dT$,
\begin{equation}\label{xid:wat:eq:Eulerladder}
 \Theta_T\mathcal I_m=\mathcal I_{m-1}\quad(m\geq2),
 \qquad
 \Theta_T\mathcal I_1=D_p(T)D_q(T).
\end{equation}
Thus every fixed integration order has a finite expression in double
polylogarithms of weight $m+2$ and depth at most two.
\end{theorem}

\begin{proof}
Expand the logarithms for $|a|,|b|<1$ and integrate term by term.
Absolute convergence gives
\[
 \mathcal L_m(a,b)
 =\sum_{j,k\geq1}\frac{a^jb^k}{jk(j+k)^m}.
\]
Use $1/(jk)=(1/j+1/k)/(j+k)$ and write $n=j+k$.
The two resulting sums are exactly those in
\eqref{xid:wat:eq:LmMPL}. Along the indicated imaginary radial paths,
$1-\lambda at$ and $1-\lambda bt$ never vanish for $0\leq t\leq1$;
analytic continuation gives the asserted extension. The real-part
calculation establishing \eqref{xid:wat:eq:Jintegral} is unchanged by
the logarithmic weight. The four-term expansion and the scaling
$t=Ts$ prove \eqref{xid:wat:eq:Imint}. Differentiating its normalized
Cauchy integration kernel gives \eqref{xid:wat:eq:Eulerladder}.
Near $t=0$ the product $D_p(t)D_q(t)$ is $O(t^2)$, which
justifies every differentiation and fixes the primitive constants.
\end{proof}

One useful diagonal identity follows without another integration:
\begin{equation}\label{xid:wat:eq:diagonalMPL}
 \mathcal J_m(A,A)
 =2^{-m-1}\operatorname{Li}_{m+1,1}(-A^2,1)
   -2\operatorname{Re}\operatorname{Li}_{m+1,1}(iA,1).
\end{equation}
Indeed, polarization of
$\log(1-at)+\log(1+at)=\log(1-a^2t^2)$ gives
\[
 \mathcal L_m(a,-a)
 =2^{-m}\operatorname{Li}_{m+1,1}(a^2,1)
   -\operatorname{Li}_{m+1,1}(a,1)
   -\operatorname{Li}_{m+1,1}(-a,1),
\]
where $t\mapsto t^2$ accounts for $2^{-m}$. Combine this with
$\mathcal L_m(a,a)=2\operatorname{Li}_{m+1,1}(a,1)$.
The arguments in \eqref{xid:wat:eq:diagonalMPL} lie off the positive-real
cut, so the same radial continuation is unambiguous.

\subsection{An exact obstruction to the one-product ansatz}

The earlier paragraph's failure of the unweighted factorization can
be made into a very short exact statement. Set
\[
 u(p)=\frac{p}{1+\sqrt{1+p^2}},\qquad u(0)=0.
\]
The unweighted integral is a function of $u(p)u(q)$. The weighted
integral has no such analytic factorization, even into an unspecified
one-variable function.

\begin{proposition}\label{xid:wat:prop:nonfactor}
There is no function $G$ analytic near zero for which
$K(p,q)=G(u(p)u(q))$ holds for all sufficiently small real $p,q$.
\end{proposition}

\begin{proof}
For $|p|,|q|<1$, integration of the absolutely convergent arctangent
series gives
\begin{equation}\label{xid:wat:eq:doubleTaylor}
 K(p,q)=\sum_{j,k\geq0}
 \frac{(-1)^{j+k}p^{2j+1}q^{2k+1}}{(2j+1)(2k+1)}
 \frac{4^{j+k}}{(2j+2k+1)\binom{2j+2k}{j+k}}.
\end{equation}
In particular,
\[
 K(p,q)=pq-\frac29(p^3q+pq^3)
       +\text{terms of total degree at least six},
\]
whereas
\[
 u(p)u(q)=\frac{pq}{4}-\frac{p^3q+pq^3}{16}
       +\text{terms of total degree at least six}.
\]
The coefficient of $pq$ would force $G'(0)=4$, and hence the
coefficient of $p^3q$ would have to be $-1/4$, contradicting $-2/9$.
Terms of order at least two in $G$ cannot contain $p^3q$.
\end{proof}

This is an obstruction only to the stated functional ansatz. It says
nothing about independence of particular constants. Theorem~\ref{xid:wat:thm:main}
shows explicitly that a short formula does exist after changing to
the hyperbolic sum and difference of the parameters.
Unlike a bounded search for one $\chi_3$, the coefficient mismatch
proves the exact assertion for every analytic $G$.

\paragraph{Reproduction and next questions.}
The accompanying verification program checks
\eqref{xid:wat:eq:Pderivative} and the nonfactorization coefficients
symbolically, compares both classical formulas with direct real
quadrature, verifies the hyperbolic mixed derivative and the
ordinary primitive at several endpoints,
tests the inverse-hyperbolic family at signed and nearly saturated
endpoints, and checks the higher integration orders against independent nested
series in a common convergence disk. These calculations audit
normalizations and branch choices; the proofs above establish the
identities. Natural next tasks are to identify additional endpoints
at which the four-kernel ordinary primitive compresses to real
one-variable kernels, to seek analogous addition formulas for higher
Euler integration orders, and to determine which specializations
of that higher-weight ladder admit classical polylogarithm reductions.
The Gaussian $S_6$ and current $S_8$
candidates remain unresolved by this argument.

% END sections/weighted_arctangent.tex


% BEGIN sections/audit_research.tex
\section{Source audit, reproducibility, and further research}
\label{xid:sec:audit}

\subsection{Corrections and exact claim status}

Two narrow issues remain in the pinned canonical
\path{chapters/03-algebraic.tex}. Both had been flagged in an earlier
continuation and are confirmed here; they are not presented as newly
discovered mistakes. The proposed patch preserves the existing labels
and the proved identities.

\paragraph{The real logarithm in \texttt{cleo:eq:lintanh}.}
The formula concerns real $|\lambda|<1$, with
$\theta=\arcsin\lambda$. Its term
$\theta\log\tan(\theta/2)$ must be
\[
 \theta\log\left|\tan\frac\theta2\right|,
\]
with continuous value zero at $\theta=0$, unless the parameter range is
restricted to positive $\lambda$. For negative $\lambda$ the displayed
logarithm otherwise introduces an imaginary term into a real integral.
For example, at $\lambda=-1/2$ the spurious principal-logarithm term is
$-i\pi^2/6$. The surrounding Clausen terms are unchanged. The real
identity is checked by differentiation in $\theta$, giving
$\theta/\sin\theta$ on both sides, and the zero value at the origin.

\paragraph{The description of the unweighted diagonal.}
The phrase ``the simplest non-elementary diagonal'' is replaced by
``the $(1,1)$ diagonal has the Legendre--chi representation.'' The exact
formula
\[
 \int_0^{\pi/2}\arctan^2(\sin x)\,dx
       =\pi\chi_2(3-2\sqrt2)
\]
is retained. A representation alone does not prove minimality or
nonreduction in a specified algebra of constants. This correction does
not concern the different, weighted integral evaluated in
\eqref{xid:wat:eq:K11}.

\medskip
\noindent The package contains a unified patch and its before/after
hashes. The patch passed a dry-run application and was applied to a
temporary copy, whose bytes matched the intended correction. The
baseline itself was unchanged.

\begin{longtable}{@{}>{\raggedright\arraybackslash}p{0.29\textwidth}>{\raggedright\arraybackslash}p{0.64\textwidth}@{}}
\toprule
Source item & Status after this report\\
\midrule
\endfirsthead
\toprule Source item & Status after this report\\\midrule
\endhead
\texttt{gauss:eq:S4-closed} & Already proved in the canonical
\path{04-S4-proof.tex}; this report does not claim a new proof of its
rational octahedral certificate.\\[3pt]
\texttt{cycloquot:conj:S6} & The short $S_6$ identity remains
conjectural. The existing incomplete-beta and reflection generators do
not supply the missing odd-weight relation.\\[3pt]
\texttt{s8new:conj:S8} & The current $S_8$ candidate remains
conjectural. It is distinct from the rejected earlier vector recorded
under \texttt{research:prop:S8-rejected}.\\[3pt]
Higher golden vectors & The numerical $\Li_6$--$\Li_9$ candidates in
\path{03-algebraic.tex} have not acquired proofs here.\\[3pt]
Colliding unequal grids & Resolved for every pair of Stieltjes indices
and argument orders by Section~\ref{xid:cld:section}, including the full
subtraction polynomial and convergent collision matching.\\[3pt]
Nonlinear local coordinates & Resolved for orientation-preserving
analytic local diffeomorphisms by Section~\ref{xid:sec:nonlinear}, including
the composition law and finite-jet dependence.\\[3pt]
Three unequal frequencies & Resolved under pairwise grid separation
by Section~\ref{xid:sec:triple-dilation}, including all argument and order
coefficients. Colliding triples remain a separate problem.\\[3pt]
Fully shifted height-one sums & An all-depth extension, with the
depth-one literature credited. The inner shift must not be confused
with the older outer-only shift or with simultaneous changes of all
exponents.\\[3pt]
Weighted arctangent family & Evaluated for all real parameters, with
ordinary and Euler primitives and a direct inverse-hyperbolic companion.
No global priority or minimality claim is made.\\
\bottomrule
\end{longtable}

The separated-grid theorem in the incoming report is valid under its
stated separation hypothesis. Extending its distributional product
directly through collision would be an invalid use of that theorem,
rather than an error in the separated statement. Similarly, omitting
$\Delta_{m,r}(a)$ from the fully shifted harmonic formula would give an
incorrect extension; the present report does not attribute that omitted
correction to an existing source. These distinctions matter when the new
sections are integrated.

\subsection{What the computational artifacts establish}

The analytic proofs establish the parameter-uniform identities. The
computer checks target signs, scales, coordinate constants, finite
coefficient extraction, and independent numerical evaluations. They do
not infer identities from small residuals. All numerical results are
ordinary floating-point audits; no result file is described as an
interval certificate.

\begin{center}
\small
\begin{tabular}{@{}>{\raggedright\arraybackslash}p{0.23\textwidth}>{\raggedright\arraybackslash}p{0.42\textwidth}>{\raggedright\arraybackslash}p{0.25\textwidth}@{}}
\toprule
Family & Independent audit & Observed residual\\
\midrule
Nonlinear coordinates & 213 exact rational-series checks, including
cutoff constants and composition; 4 subtracted integrals at 70 digits
& Absolute $<4.0\cdot10^{-36}$\\[4pt]
Colliding products & 33 exact coefficient checks; 9 polygamma integrals
at 44 digits and 2 Stieltjes derivative integrals at 36 digits
& Scaled $<6.7\cdot10^{-45}$ and $<2.8\cdot10^{-31}$\\[4pt]
Fully shifted sums & 72 checks at 75 digits, using Gamma-ratio
Dirichlet tails and direct local differentiation
& Absolute $<8.5\cdot10^{-52}$\\[4pt]
Three frequencies & 2 exact rational Bernoulli integrals compared with
12 weighted-kernel Mellin integrals at 60 digits
& Absolute $<7.2\cdot10^{-66}$\\[4pt]
Weighted primitives & 78 comparisons at 80 digits; symbolic checks of
the mixed derivatives, primitive, and coefficient obstruction
& Absolute $<1.3\cdot10^{-79}$\\
\bottomrule
\end{tabular}
\end{center}

There are 167 numerical comparisons in total. The difference between
working precision and observed residual is expected: some checks use
finite Taylor or asymptotic tails, while others compare small exact
values with direct quadrature. The results record the actual precision
and method rather than a uniform advertised number of correct digits.
For the harmonic continuation, two larger-tail repeats give residuals
below $1.7\cdot10^{-64}$, which tests the identified truncation risk.

The seven mathematical scripts are independent of the TeX build. They
use Python, SymPy, and mpmath. A driver reruns them, writes their logs,
and validates the result files. The exact collision generator also
accepts $m,n,r,k$ as command-line arguments and prints the corresponding
finite symbolic identity, with formal symbols for ordinary zeta
derivatives. The package includes the source manifest, a claim ledger,
integration instructions, the correction patch, and both modular and
standalone TeX. No numerical relation search is needed to reproduce any
proof in this report.

\subsection{Proposed exact research problems}
\label{xid:sec:further-research}

The following questions concern identities and structural reductions.
They are questions, not additional theorems. Each starts from a specific
formula proved above.

\begin{enumerate}[label=\textbf{Q\arabic*.},leftmargin=*]
\item \textbf{Mixed collisions of three unequal grids.}
Remove pairwise separation in Theorem~\ref{xid:thm:triple-FP}. Classify
pairwise collision points and triple collision points, write their
complete local subtraction polynomials, and determine whether the
regular constant in every simultaneous collision limit agrees with the
merged ambient finite part. The two-grid result suggests a positive
matching statement for suitably specified collision coordinates, but
multiple rates can introduce distinct logarithmic terms and must be
tracked explicitly.

\item \textbf{An all-factor finite congruence formula.}
For $\ell\ge4$ pairwise separated integer grids, generalize the six
sign-cone proof to a finite collection of colored multiple Tornheim or
Mordell--Tornheim kernels on the rank-$(\ell-1)$ lattice
$\sum p_jn_j=0$. Determine a useful minimal finite character formula,
then give every ambient-coordinate Stieltjes coefficient with complete
local numerators. The challenge is an exact finite organization of the
lattice, rather than existence of meromorphic continuation alone.

\item \textbf{Which odd-order cancellations survive at more than two
factors?}
Corollary~\ref{xid:cld:diagonal-collapse} cancels an entire antisymmetric
two-factor spectral kernel. Identify a permutation or character
projection for three or more factors that removes the genuinely
multiple spectral terms and leaves only local Stieltjes data. A useful
first case has equal Stieltjes indices and one total argument derivative.

\item \textbf{Nonlinear transport of merged products.}
Theorem~\ref{xid:thm:coordinate-monomial} applies to the full singular
expansion of a colliding product. Compose it with
\eqref{xid:cld:common-subtraction} to obtain a single closed multivariate
generator for every nonlinear pullback of every two-grid product.
Determine sharp vanishing thresholds in the two Stieltjes indices.
For a nonunit tangent, factor out the affine contribution before
classifying the curvature terms.

\item \textbf{A finite algebra for coordinate composition.}
Express the coefficients of \eqref{xid:eq:C-generator} and
\eqref{xid:eq:delta-pullback} in a consistently normalized Bell-polynomial
basis and identify the resulting composition matrices. Which
logarithmic singular germs have zero defect for every coordinate tangent
to the identity? The theorem supplies a direct coefficient test, but a
classification by singular order and logarithmic degree could expose
additional exact cancellations.

\item \textbf{Higher regular Laurent coefficients of the harmonic
entire difference.}
Evaluate $\partial_s^kD(-m,u;a)$ for $k\ge1$. Its value at $k=0$
is a finite Stirling sum, whereas differentiation introduces logarithmic
Mellin moments of the analytic Gauss block. Seek closed expressions in
Gamma derivatives, generalized Stieltjes constants, and a precisely
specified class of multiple polylogarithms. Do not assume the rational
polynomial closure at $k=0$ persists for $k\ge1$.

\item \textbf{Identify the Stirling polynomial system exactly.}
Compare
$\sum_j\left\{\begin{smallmatrix}m+1\\j\end{smallmatrix}\right\}
(a-1)^{\underline j}/j^{r+1}$
with poly-Bernoulli and related polynomial systems, including the
falling-factorial convention and the shift by one. An exact generating
function identification could yield addition, reflection, and
differentiation identities for the complete harmonic correction.

\item \textbf{Independent inner exponents.}
Replace the inner exponent string $(1,\ldots,1)$ in $M_r$ by independent
parameters, keeping the common shift. Determine the smallest generating
kernel whose singular part still separates by a Gamma factor, and
identify the extra entire corrections. This would connect the present
height-one calculus to the incoming all-exponent nested jets without
conflating their spectral directions.

\item \textbf{Rational shifts and character projections.}
Apply finite Dirichlet-character sums to the $a$-dependence of the
Laurent formulas. Determine exact reductions of their Gamma-polygamma
coefficients to Dirichlet $L$-values and their derivatives, and identify
which character parities eliminate the polynomial Stirling correction.
The finite correction should be summed before imposing parity.

\item \textbf{A real addition law for higher Euler primitives.}
The first weighted arctangent integral has the two-term formula
\eqref{xid:wat:eq:main}; the higher Euler primitives currently have the
finite depth-two expression \eqref{xid:wat:eq:Imclosed}. Determine which
higher integration orders admit a comparable addition formula in
$\alpha\pm\beta$, and which require irreducible two-variable functions.
The known exact coefficient obstruction concerns one product variable
only; it does not rule out these other organizations.

\item \textbf{Relations among the algebraic endpoint values.}
The weighted Cleo identity combines logarithms of the silver and golden
units with $\chi_2$ and $\chi_3$ at explicit algebraic arguments. Seek
functional-equation certificates for further reductions at these
arguments. A target basis should be specified in advance, and failure
of a finite relation search should not be interpreted as independence.

\item \textbf{The short $S_6$ and current $S_8$ candidates.}
Extend the convergent octahedral or lifted functional-equation mechanism
that proved $S_4$ to weights seven and nine. The existing reflection
projection cancels at the relevant odd total weights, so a successful
proof needs an additional functional relation. A useful deliverable
would be an exact rational row certificate with an independently
checked analytic origin for every row. The frozen candidate vectors,
not newly fitted alternatives, should remain the targets.

\item \textbf{Formalization of the local finite-part lemmas.}
Formalize the elementary Taylor-subtraction coefficient lemma first,
including its exact test-function and coordinate conventions. It can
then serve as a common foundation for nonlinear contacts and colliding
products. The finite rational coefficient computations are suitable
for checked algebraic certificates, while the analytic continuation
and dominated-differentiation steps remain separate proof obligations.
\end{enumerate}

The most immediate continuations are Q1, Q4, and Q6: their inputs are
explicitly available in this report, and each asks for an exact identity
whose first unresolved term can be computed without a numerical basis
guess. The $S_6$ and $S_8$ problems remain more structurally demanding;
the present endpoint calculus supplies normalization discipline, not
the missing weight-seven or weight-nine functional equation.

% END sections/audit_research.tex

\clearpage

% BEGIN references.tex
\begin{thebibliography}{99}
\small
\raggedright
\setlength{\itemsep}{3pt}
\setlength{\parskip}{0pt}
\addcontentsline{toc}{section}{References}

\bibitem{ProveIt}
V. Reshetnikov and contributors,
\emph{ProveIt: Polylogarithms unified manuscript},
repository snapshot \texttt{445f754610e1377939235794de84ed9075fc09f5},
10 October 2026.
\url{https://github.com/VladimirReshetnikov/ProveIt/tree/445f754610e1377939235794de84ed9075fc09f5/Analysis/Polylogarithms/docs/manuscript}.

\bibitem{IncomingGauss}
\emph{ProveIt Gauss--Hurwitz continuation}, incoming research archive
\path{ProveIt_Gauss_Hurwitz_2026-10-10.zip}, in the pinned
\path{docs/incoming} snapshot of \cite{ProveIt}.

\bibitem{IncomingMellin}
\emph{ProveIt Mellin--Dilation continuation}, incoming research archive
\path{ProveIt_Mellin_Dilation_2026-10-10.zip}, in the pinned snapshot.
In particular, its audit and research section asks for colliding
unequal dilations and nonlinear coordinates.

\bibitem{IncomingNested}
\emph{ProveIt Nested Harmonic Jets}, incoming research archive
\path{ProveIt_Nested_Harmonic_Jets_2026-10-10.zip}, in the pinned snapshot.

\bibitem{IncomingContinuation}
\emph{ProveIt Polylogarithm--Stieltjes Continuation}, incoming archive
\path{ProveIt_Polylogarithm_Stieltjes_Continuation_2026-10-10.zip},
in the pinned snapshot.

\bibitem{IncomingIdentities}
\emph{ProveIt Stieltjes--Harmonic Identities}, incoming research archive
\path{ProveIt_Stieltjes_Harmonic_Identities_2026-10-10.zip},
in the pinned snapshot. In particular, its further-research section
asks for unequal triple frequencies and collision matching.

\bibitem{EspinosaMoll1}
O. Espinosa and V. H. Moll,
\emph{On Some Integrals Involving the Hurwitz Zeta Function: Part 1},
Ramanujan Journal \textbf{6} (2002), 159--188.
\url{https://arxiv.org/abs/math/0012078}.
Theorem 3.1 is the classical coincident two-factor spectral input.

\bibitem{EspinosaMoll2}
O. Espinosa and V. H. Moll,
\emph{On Some Integrals Involving the Hurwitz Zeta Function: Part 2},
Ramanujan Journal \textbf{6} (2002), 449--468.
\url{https://arxiv.org/abs/math/0107082}.

\bibitem{ZhaoColored}
J. Zhao, \emph{A note on colored Tornheim's double series},
Integers \textbf{10} (2010), A59, 879--882.
\url{https://arxiv.org/abs/0907.5106}.

\bibitem{hsh:Kargin2023}
L. Karg\i n, A. Dil, M. Cenkci and M. Can,
\emph{On the Stieltjes constants with respect to harmonic zeta functions},
Journal of Mathematical Analysis and Applications (2023), article 127302.
\url{https://doi.org/10.1016/j.jmaa.2023.127302};
\url{https://arxiv.org/abs/2304.03517}.

\bibitem{hsh:KaraCan2026}
M. Kara \"Ozt\"urk and M. Can,
\emph{Series involving parametric harmonic zeta function},
arXiv:2606.27827, posted 26 June 2026.
\url{https://arxiv.org/abs/2606.27827}.
The primary abstract was retrieved; full-text access was unavailable
in this research environment. Used for scope and attribution, not as
an unexamined proof dependency.

\bibitem{DLMFhurwitz}
NIST Digital Library of Mathematical Functions,
\S25.11, \emph{Hurwitz Zeta Function}.
\url{https://dlmf.nist.gov/25.11}.

\bibitem{hsh:DLMFPolygamma}
NIST Digital Library of Mathematical Functions,
\S5.15, \emph{Polygamma Functions}.
\url{https://dlmf.nist.gov/5.15}.

\bibitem{hsh:DLMFGamma}
NIST Digital Library of Mathematical Functions,
\S5.11, \emph{Asymptotic Expansions}.
\url{https://dlmf.nist.gov/5.11}.

\bibitem{hsh:DLMFGauss}
NIST Digital Library of Mathematical Functions,
\S15.10, \emph{Hypergeometric Differential Equation}, especially
the connection formula 15.10.21.
\url{https://dlmf.nist.gov/15.10#E21}.

\bibitem{hsh:DLMFBernoulli}
NIST Digital Library of Mathematical Functions,
\S24.2, \emph{Definitions and Generating Functions}.
\url{https://dlmf.nist.gov/24.2}.

\bibitem{wat:DLMF}
NIST Digital Library of Mathematical Functions,
\S25.12, \emph{Polylogarithms}.
\url{https://dlmf.nist.gov/25.12}.

\bibitem{wat:DavidH}
David H., \emph{Simplifying a certain polylogarithmic sum in two variables},
Mathematics Stack Exchange, question posted 17 September 2015,
self-answer posted 1 May 2016.
\url{https://math.stackexchange.com/questions/1439117/}.

\bibitem{wat:Lee}
Sangchul Lee, answer to
\emph{Integral $\int_0^1\arctan^2x/\sqrt{1-x^2}\,dx$},
Mathematics Stack Exchange, 11 November 2013.
\url{https://math.stackexchange.com/questions/555882/}.

\end{thebibliography}

% END references.tex

\end{document}
